\documentclass[11pt,reqno]{amsart}

\usepackage[margin=1.1in]{geometry}
\usepackage{amsmath,amssymb,mathtools,amsthm,stmaryrd,mathrsfs}
\usepackage{booktabs,array,microtype,aliascnt}
\usepackage[hidelinks]{hyperref}
\usepackage[capitalize,noabbrev]{cleveref}

\numberwithin{equation}{section}
\setcounter{MaxMatrixCols}{12}
\allowdisplaybreaks

\newtheorem{theorem}{Theorem}[section]
\newaliascnt{proposition}{theorem}
\newtheorem{proposition}[proposition]{Proposition}
\aliascntresetthe{proposition}
\newaliascnt{lemma}{theorem}
\newtheorem{lemma}[lemma]{Lemma}
\aliascntresetthe{lemma}
\newaliascnt{corollary}{theorem}
\newtheorem{corollary}[corollary]{Corollary}
\aliascntresetthe{corollary}
\newaliascnt{definition}{theorem}
\newtheorem{definition}[definition]{Definition}
\aliascntresetthe{definition}
\theoremstyle{remark}
\newaliascnt{remark}{theorem}
\newtheorem{remark}[remark]{Remark}
\aliascntresetthe{remark}
\newtheorem*{maintheorem}{Main Theorem}

\crefname{theorem}{Theorem}{Theorems}
\Crefname{theorem}{Theorem}{Theorems}
\crefname{proposition}{Proposition}{Propositions}
\Crefname{proposition}{Proposition}{Propositions}
\crefname{lemma}{Lemma}{Lemmas}
\Crefname{lemma}{Lemma}{Lemmas}
\crefname{corollary}{Corollary}{Corollaries}
\Crefname{corollary}{Corollary}{Corollaries}
\crefname{definition}{Definition}{Definitions}
\Crefname{definition}{Definition}{Definitions}
\crefname{remark}{Remark}{Remarks}
\Crefname{remark}{Remark}{Remarks}

\newcommand{\kk}{\Bbbk}
\newcommand{\NN}{\mathbb{N}}
\newcommand{\ZZ}{\mathbb{Z}}
\newcommand{\RR}{\mathbb{R}}
\newcommand{\CC}{\mathbb{C}}
\newcommand{\cA}{\mathcal{A}}
\newcommand{\cC}{\mathcal{C}}
\newcommand{\cF}{\mathcal{F}}
\newcommand{\cL}{\mathcal{L}}
\newcommand{\cS}{\mathcal{S}}
\newcommand{\cU}{\mathcal{U}}
\newcommand{\cX}{\mathcal{X}}
\newcommand{\SR}{\mathscr S}
\newcommand{\UR}{\mathscr U}
\newcommand{\As}{\mathcal A_s}
\newcommand{\Asup}{\mathcal A_s^{\mathrm{up}}}
\newcommand{\Uneg}{\mathcal U_u^{-}}
\newcommand{\Mu}{\mathcal M_u}
\newcommand{\Muwt}[1]{\mathcal M_{u,#1}}
\newcommand{\Ugen}{\mathscr U_{\mathrm{gen}}}
\newcommand{\udisc}{u_{\Delta}}
\newcommand{\g}{{\sf g}}
\newcommand{\wt}{\operatorname{wt}}
\newcommand{\Frac}{\operatorname{Frac}}
\newcommand{\Spec}{\operatorname{Spec}}
\newcommand{\Sym}{\operatorname{Sym}}
\newcommand{\GL}{\operatorname{GL}}
\newcommand{\SL}{\operatorname{SL}}
\newcommand{\Rep}{\operatorname{Rep}}
\newcommand{\SI}{\operatorname{SI}}
\newcommand{\ord}{\operatorname{ord}}
\newcommand{\initer}{\operatorname{in}}
\newcommand{\dlog}{\,d\!\log}
\newcommand{\pos}[1]{\left[#1\right]_{+}}
\newcommand{\midalg}{\Mu}
\newcommand{\Schur}{\mathbf{S}}
\newcommand{\Specht}{\mathsf{S}}
\newcommand{\bmat}[1]{\begin{bmatrix}#1\end{bmatrix}}

\title{A polyhedral formula for $n\times2\times2$ Kronecker coefficients via cluster algebras}
\author{Jiarui Fei}
\address{School of Mathematical Sciences, Shanghai Jiao Tong University, Shanghai, China}
\email{jiarui@sjtu.edu.cn}
\author{Chenxin Xue}
\address{School of Mathematical Sciences, Shanghai Jiao Tong University, Shanghai, China}
\email{xcx0711@sjtu.edu.cn}
\thanks{The authors were supported in part by the National Natural Science Foundation of China (Nos.~12131015 and 12571038).}
\subjclass[2020]{Primary 13F60, 05E10; Secondary 13A50, 20C30}
\keywords{Kronecker coefficient, tensor invariant, cluster algebra, theta function, logarithmic volume form}
\date{}

\begin{document}

\begin{abstract}
We study the triple-invariant algebra
\[ \UR=\kk[\kk^3\otimes\kk^2\otimes\kk^2]^{U_3\times U_2\times U_2}. \]
A quotient slice and the induced logarithmic top form determine a signed
Markov chart, realized as the fiber $\zeta=-1$ of an ordinary cluster
family.  We prove
\[ \UR=\Uneg=\Mu[\udisc], \]
where $\Mu$ is a specialized middle cluster algebra and $\udisc$ is the
discriminant of weight $(220;22;22)$.  Its theta cone has a sixteen-element
Hilbert basis.  Pairing its positive- and negative-degree generators reduces each
triple-weight space to a single discriminant level determined by the weight.
Counting the resulting two-dimensional slice gives an explicit nonnegative finite-sum formula.
Determinant reduction extends the formula to all $n\times2\times2$
Kronecker coefficients.
\end{abstract}

\maketitle

\section*{Introduction}

Kronecker coefficients are the structure constants for tensor products of
irreducible representations of symmetric groups.  Despite their elementary
definition, no general positive combinatorial rule is known.  The case in
which two indexing partitions have at most two rows is exceptional.  General
treatments of this case include the explicit formulas of
Remmel--Whitehead and Rosas \cite{RW,Rosas}, the geometric approach of
Adsul--Subrahmanyam \cite{AS}, the crystal basis and positive formula of
Blasiak--Mulmuley--Sohoni \cite{BMS}, the chamber quasi-polynomials of
Briand--Orellana--Rosas \cite{BOR09}, and the vector-partition formulation of
Mishna--Rosas--Sundaram \cite{MRS}.  These works already provide effective
formulas.

Our point is not primarily to compute these already-understood
coefficients faster.  Rather, we give a geometric and canonical-basis
explanation of why a positive polyhedral formula exists.  Concretely, we
derive an explicit nonnegative finite-sum formula by realizing the coefficients in a
coefficient-specialized Markov cluster algebra.  For each triple weight, the
relevant theta indices (also known as $\g$-vectors) lie on a single slice of a
rational cone, and counting that slice gives the formula.  This description
also yields piecewise quasi-polynomial behavior with quasiperiod dividing
two; the parity bound is sharp.  Determinant extension then covers all
$n\times2\times2$ formats.

The basic invariant algebra is
\[ \UR=\kk[\kk^3\otimes\kk^2\otimes\kk^2]^{U_3\times U_2\times U_2}. \]
Its triple-weight space of weight $(\lambda;\mu;\nu)$ has dimension
$g_{\lambda,\mu,\nu}$.  Determinant extension identifies the stable
$4\times2\times2$ algebra with a polynomial extension of $\UR$, and
rectangular translation then reduces every $n\times2\times2$ coefficient
to the $(3,2,2)$ calculation.

There are two obstacles to a direct cluster treatment.  The known ordinary
cluster structure lies on the double-invariant algebra
\[ \SR =\kk[\kk^3\otimes\kk^2\otimes\kk^2]^{U_3\times U_2}, \]
whereas $\UR=\SR^{U_2}$ is the kernel of one additional locally nilpotent
root derivation.  Moreover, the natural Markov exchange on the affine
quotient $Y_{322}=\Spec\UR$ has a negative sign.  We avoid introducing a
separate Laurent-phenomenon formalism: an auxiliary coefficient $\zeta$
embeds the signed exchanges in an ordinary cluster family over
$\kk[\zeta^{\pm1}]$, and the model on $Y_{322}$ is recovered on the fiber
$\zeta=-1$.

We construct a birational chart on $Y_{322}$ from a local slice with
rational coordinates $q_1,\ldots,q_7$.  Contracting the logarithmic top form
of the double seed
with the root derivation gives
\[ \omega_{322} =-\frac{dq_1\wedge\cdots\wedge dq_7}{q_1q_2^3}. \]
Comparing a proposed torus chart with this form tests algebraic independence
and determines the anticanonical boundary. This comparison gives the two
adjacent logarithmic charts.
%\[ (u_1,u_5,u_7;u_2,u_3,u_4,u_{13}), \qquad (u_1,u_5,u_8;u_2,u_3,u_4,u_{13}), \]
Their rational transition gives the exchange relation
%\[ u_7u_8=u_4u_5^2+u_1^2u_{13}. \]
and the Jacobian constants record the corresponding orientation signs.
The resulting model turns out to be a signed Markov cluster algebra on the negative coefficient fiber.

The four-frozen chart also comes from the trivalent maximal parabolic of type
$D_5$.  Its degree $-1$ representation is
$\kk^3\otimes\kk^2\otimes\kk^2$.  The seven seed functions are restrictions
of orthogonal Pl\"ucker covariants, and the exchange matrix is an integral
Gale dual of their $D_5$-weights.  Appendix~\ref{app:D5-parabolic} gives the
details; this interpretation is not used in the proofs.

Before using theta functions, we prove that thirteen normalized
highest-weight invariants generate $\UR$.  Root-string calculations and two
rational slices identify the invariant ring after localization at $u_1$ and
at $u_2$.  A regular-sequence computation and an elementary
intersection argument then give the global equality.  Specializing the initial
and three adjacent Laurent rings of the $\zeta$-family defines a four-chart algebra
$\Uneg$; a direct UFD denominator argument proves
$\Uneg\subseteq\UR$.

The mutable quiver is the Markov quiver, hence the surface quiver of the
once-punctured torus.  For theta functions we use Zhou's full Fock--Goncharov
theorem for the principal-coefficient Markov seed together with valuative
independence under coefficient specialization \cite{Zhou,CMMM}.  If
$\Sigma_{g,p}$ is triangulable and $p\ge2$, Li constructs an open
immersion from the cluster variety into the spectrum of the tagged skein
algebra with complement of codimension at least two, and deduces
$\operatorname{Sk}^{\mathrm{ta}}(\Sigma)=\mathcal U(\Sigma)$, equivalently
$\operatorname{mid}(\mathcal A)=\operatorname{up}(\mathcal A)$
\cite[Theorem~1.1(1) and Corollary~1.2]{Li}.  For $p=1$, Li proves the
analogous codimension-two statement only for an enlarged cluster variety
\cite[Theorem~1.1(2)]{Li}.  Neither result identifies the middle and upper
algebras for the ordinary once-punctured-torus seed, and neither treats the
non-inverted boundary variables or the specialization $\zeta=-1$ used
here.

After specialization, regularity along the four frozen boundary divisors is
encoded by a rational cone $\Xi\subset\RR^7$.  Its sixteen primitive
Hilbert-basis labels generate a specialized middle cluster algebra $\Mu$.
Twelve of the thirteen invariant generators lie in $\Mu$, while the
remaining one is the discriminant
\[ \udisc=u_{220,22}^{22}. \]
Two identities place $\udisc$ in all four Laurent charts.  Together with
the preceding inclusions, they give
\[ \UR=\Uneg=\Mu[\udisc]. \]
%For orientation, the main objects are related by
%\[ \mathcal A_{\zeta} \rightsquigarrow \mathcal A_{\zeta=-1}, \qquad
%   \Mu \subseteq \Uneg=\cL_\circ^-\cap \cL_1^-\cap \cL_5^-\cap \cL_7^- \subseteq \UR, \qquad
%   \UR=\Mu[\udisc]. \]
%Here $\mathcal A_{\zeta}$ is the ordinary coefficient family,
%$\mathcal A_{\zeta=-1}$ is its negative fiber, and $\Uneg$ is the
%intersection of the four Laurent charts on that fiber.
The four negative-degree Hilbert generators are, up to nonzero scalars,
$\udisc$ times the four positive-degree generators.  This pairing reduces
each fixed triple-weight space to the discriminant level determined by its
weight.  The cone slice is two-dimensional; eliminating its
linear constraints gives the finite sum in \cref{thm:closed-formula}.
Appendix~\ref{app:folded-DT} records the folded Donaldson--Thomas
transformation behind the pairing; it is not needed in the proof.  Two further
coefficient-fiber models are recorded without proofs in
Appendix~\ref{app:other-models}.

\begin{maintheorem}
For
\[ \UR=\kk[\kk^3\otimes\kk^2\otimes\kk^2]^{U_3\times U_2\times U_2}, \]
the following statements hold.
\begin{enumerate}
\item The thirteen invariants in \eqref{eq:thirteen-invariant-weights}
generate $\UR$, and
\[ \UR=\Uneg=\Mu[\udisc]. \]
\item The specialized middle algebra $\Mu$ has theta basis
\[ \{\theta_v\mid v\in\Xi\cap\ZZ^7\}, \]
the sixteen labels in \eqref{eq:Xi-Hilbert-basis} form the Hilbert
basis of $\Xi\cap\ZZ^7$, and the corresponding theta functions form a Khovanskii generating set for $\Mu$.
\item For an admissible triple weight
$\chi=(\lambda;\mu;\nu)$, let $\ell_0$ be defined by
\eqref{eq:A-and-level}, and let
$\chi_\Delta=\wt(\udisc)$.  Then
\[ g_{\lambda,\mu,\nu} =\#\{v\in\Xi\cap\ZZ^7\mid vW=\chi-\ell_0\chi_\Delta\}, \]
and the count is the nonnegative finite sum in
\eqref{eq:closed-Kronecker-sum}.  Rectangular translation and stabilization
extend the same formula to every $n\times2\times2$ format.
\end{enumerate}
\end{maintheorem}

Every finite calculation used in the paper is small enough to be checked by
hand from the matrices and formulas.  The ancillary file
\path{Kron322.py} is provided only to shorten the exposition and to make these
checks convenient for readers to reproduce; it is not used in place of a
mathematical argument.

The paper is organized as follows.  \Cref{sec:invariants} recalls the
invariant-theoretic realization of Kronecker coefficients and determinant
reduction.  \Cref{sec:cluster-background} fixes the cluster, upper, and theta
conventions.  The quotient slice, logarithmic form, and ordinary coefficient
family are constructed in \cref{sec:slice-volume}.  The ring equalities and
the specialized middle algebra are proved in \cref{sec:ring-equalities}; the
reduction to one theta level and the closed formula occupy
\cref{sec:polyhedral}.  Appendix~A recalls the upstairs cluster algebra,
Appendix~B gives the regular-sequence computation, Appendix~C gives the folded
DT interpretation, Appendix~D explains the $D_5$-parabolic origin of the
four-frozen seed, and Appendix~E records two additional coefficient-fiber
models without proofs.  The ancillary script
\path{Kron322.py}, with its README and pinned SymPy version,
performs the exact fan, Hilbert-basis, $D_5$-lattice, weight-semigroup,
Laurent, regular-sequence, folded-DT, and finite-sum checks.  Its default run
also compares the
formula with the symmetric-group character formula for all relevant triples
of size at most twelve.

Throughout, $\kk$ is an algebraically closed field of characteristic zero.
All partitions are padded by trailing zeros when necessary.  We use the
compact notation $110=(1,1,0)$, $21=(2,1)$, and so forth.

\section{Tensor invariants and Kronecker coefficients}
\label{sec:invariants}

\subsection{Highest-weight algebras}

Let $\lambda\vdash N$ be a partition.  We write $\Specht^\lambda$ for
the corresponding Specht module of the symmetric group $\mathfrak S_N$,
and $\Schur_\lambda V$ for the corresponding Schur module.  For three
partitions of $N$, the Kronecker coefficient is defined by
\[ \Specht^\mu\otimes\Specht^\nu \cong\bigoplus_{\lambda\vdash N} g_{\lambda,\mu,\nu}\Specht^\lambda. \]

For positive integers $a,b,c$, set
\[ T_{a,b,c}=\kk^a\otimes\kk^b\otimes\kk^c. \]
Throughout, $\kk[T_{a,b,c}]$ denotes the polynomial coordinate ring
$\Sym(T_{a,b,c}^{\vee})$, with the three general linear groups acting on
functions contragrediently.  We use the standard polynomial highest-weight
labels in the resulting Cauchy decomposition and suppress dual symbols on
Schur functors.  Let $U_d\subset\GL_d$ be the standard upper unitriangular
subgroup.
We use the two highest-weight algebras
\[ \SR_{a,b,c} :=\kk[T_{a,b,c}]^{U_a\times U_b}, \qquad \UR_{a,b,c} :=\kk[T_{a,b,c}]^{U_a\times U_b\times U_c}. \]
They are graded by triples of dominant weights.  When
$(a,b,c)=(3,2,2)$, we abbreviate them to $\SR$ and $\UR$.  In this
case the last factor $U_2$ is isomorphic to $\mathbb G_a$.  Set
\[ X_{\mathrm{dbl}}:=\Spec\SR, \qquad Y_{322}:=\Spec\UR, \qquad \pi:X_{\mathrm{dbl}}\longrightarrow Y_{322}. \]
Then $Y_{322}=X_{\mathrm{dbl}}/\!/\mathbb G_a$ is the affine categorical
quotient: every $\mathbb G_a$-invariant morphism from
$X_{\mathrm{dbl}}$ to an affine scheme factors uniquely through $\pi$.
All subsequent references to the quotient mean $Y_{322}$.  The function
field calculation in \cref{prop:birational-quotient-chart} gives dimension
seven, and \cref{thm:thirteen-generate} proves finite generation; hence
$Y_{322}$ is an affine variety.  Before those results, the same notation is
understood scheme-theoretically.  The standard Cauchy--Schur--Weyl
realization used below is recalled, for example, in
\cite[Chapter~I]{Macdonald}.

\begin{lemma}
\label{lem:kronecker-weight-space}
For partitions $\lambda,\mu,\nu$ of the same size satisfying
$\ell(\lambda)\le a$, $\ell(\mu)\le b$, and $\ell(\nu)\le c$, one has
\[ \dim (\UR_{a,b,c})_{\lambda,\mu,\nu} =g_{\lambda,\mu,\nu}. \]
\end{lemma}

\begin{proof}
The Cauchy decomposition and the restriction of a Schur functor to a tensor
product give
\[
 \Sym^N(\kk^a\otimes\kk^b\otimes\kk^c)
 \cong
 \bigoplus_{\lambda,\mu,\nu\vdash N}
 g_{\lambda,\mu,\nu}
 \bigl(\Schur_\lambda\kk^a\otimes
       \Schur_\mu\kk^b\otimes
       \Schur_\nu\kk^c\bigr).
\]
Each irreducible polynomial $\GL$-module has a one-dimensional highest
weight line.  Taking invariants under the three maximal unipotent subgroups
therefore retains precisely $g_{\lambda,\mu,\nu}$ copies of the indicated
triple highest-weight line.
\end{proof}

\subsection{Determinant reduction and stabilization}

Fix $b,c\ge1$, put $D=bc$, and assume $D\ge2$.  Let
$A=\kk^D$, $B=\kk^b$, $C=\kk^c$, and $E=B\otimes C$.  A tensor in
$A\otimes B\otimes C=A\otimes E$ is a square $D\times D$ matrix.  Let
\[ \Delta_{b,c} =\det\bigl(\operatorname{Flat}_{A\mid B\otimes C}\bigr) \in\Sym^D(A\otimes B\otimes C) \]
be the determinant of this flattening.

\begin{proposition}
\label{prop:det-extension}
Under the standard inclusion $\kk^{D-1}\subset\kk^D$, one has a direct
sum decomposition
\[ \UR_{D,b,c} =\bigoplus_{r\ge0}\Delta_{b,c}^{\,r}\UR_{D-1,b,c}. \]
Consequently,
\[ \UR_{bc,b,c}\cong\UR_{bc-1,b,c}\otimes_\kk\kk[\Delta_{b,c}] \]
as multigraded algebras, and
\[ \deg\Delta_{b,c} =\bigl((1^{bc}),(c^b),(b^c)\bigr). \]
\end{proposition}

\begin{proof}
The Cauchy decomposition is
\[ \Sym(A\otimes E) =\bigoplus_{\ell(\lambda)\le D} \Schur_\lambda A\otimes\Schur_\lambda E. \]
If $\lambda_D=0$, the $U_D$-highest-weight line in
$\Schur_\lambda A$ is the image of the $U_{D-1}$-highest-weight line in
$\Schur_\lambda\kk^{D-1}$.  Hence the sum of the components with
$\lambda_D=0$ is naturally $\UR_{D-1,b,c}$.

The element $\Delta_{b,c}$ spans $\det(A)\otimes\det(E)$, and
\[ \det(E)\cong(\det B)^c\otimes(\det C)^b. \]
For $\lambda=(\lambda_1,\ldots,\lambda_D)$, put $r=\lambda_D$ and
$\bar\lambda=\lambda-r(1^D)$.  The determinant-twist identities
\[ \Schur_\lambda A\cong(\det A)^r\otimes\Schur_{\bar\lambda}A, \qquad \Schur_\lambda E\cong(\det E)^r\otimes\Schur_{\bar\lambda}E \]
are realized by multiplication by $\Delta_{b,c}^{\,r}$.  After
restriction from $\GL(E)$ to $\GL(B)\times\GL(C)$, the last factor shifts
the other two highest weights by $rc(1^b)$ and $rb(1^c)$, respectively.
Thus every homogeneous component is uniquely a power of $\Delta_{b,c}$
times a component with last part zero.  The sum is direct because the power
is recovered from $\lambda_D$.
\end{proof}

\begin{corollary}
\label{cor:rectangular-translation}
Let $r=\lambda_{bc}$.  If $g_{\lambda,\mu,\nu}\ne0$, then
\[ \mu_b\ge rc,\qquad \nu_c\ge rb, \]
and
\[ g_{\lambda,\mu,\nu} =g_{\lambda-r(1^{bc}),\,\mu-rc(1^b),\,\nu-rb(1^c)}. \]
If either inequality fails, then $g_{\lambda,\mu,\nu}=0$.
\end{corollary}

\begin{proof}
Take dimensions in the homogeneous-component isomorphisms in the proof of
\cref{prop:det-extension} and apply \cref{lem:kronecker-weight-space}.
\end{proof}

\begin{corollary}
\label{cor:stable-first-factor}
For every $a\ge bc$, the standard inclusion $\kk^{bc}\subset\kk^a$
induces an isomorphism
\[ \UR_{bc,b,c}\xrightarrow{\sim}\UR_{a,b,c}, \]
where the first partition is padded by $a-bc$ zeros.  Hence
\[ \UR_{a,b,c}\cong\UR_{bc-1,b,c}\otimes_\kk\kk[x] \qquad(a\ge bc). \]
For $a>bc$, the new generator is the highest-weight $bc\times bc$
maximal minor of the $a\times bc$ flattening.
\end{corollary}

\begin{proof}
Only partitions of length at most $bc=\dim E$ occur in the Cauchy
decomposition of $\Sym(\kk^a\otimes E)$.  Their $U_a$-highest-weight
lines are already contained in the standard copy of $\kk^{bc}$.  The first
assertion follows componentwise, and the second follows from
\cref{prop:det-extension}.
\end{proof}

\begin{remark}
At the level of Kronecker coefficients, \cref{cor:rectangular-translation}
is the rectangular translation symmetry of Briand--Orellana--Rosas
\cite[Proposition~7]{BOR}; the $(4,2,2)$ case also appears in
\cite[Corollary~5.1]{Vallejo}.  The proposition above is the compatible
algebra-level statement.
\end{remark}

\section{Cluster, upper, and middle algebras}
\label{sec:cluster-background}

\subsection{Cluster and upper cluster algebras}

Let $\mu,f\ge0$, put $m=\mu+f$, let
$\cF=\kk(x_1,\ldots,x_m)$, and let
$B=(b_{ki})\in M_{\mu\times(\mu+f)}(\ZZ)$ have skew-symmetrizable
principal part.  Throughout the paper, rows are indexed by the $\mu$
mutable variables and columns by all $\mu+f$ mutable and frozen variables.
A seed is a pair $\Sigma=(\mathbf{x},B)$, where
$\mathbf{x}=(x_1,\ldots,x_m)$ is a transcendence basis of $\cF$.  The
first $\mu$ variables are mutable and the remaining $f$ variables are
frozen.  Mutation in direction $k\le\mu$ replaces $x_k$ by $x_k'$, where
\[ x_kx_k' =\prod_{b_{ki}>0}x_i^{b_{ki}} +\prod_{b_{ki}<0}x_i^{-b_{ki}}, \]
and mutates $B$ by the usual matrix rule in this convention.  We write
\[ \cL(\mathbf{x}) =\kk[x_1^{\pm1},\ldots,x_\mu^{\pm1},x_{\mu+1},\ldots,x_{\mu+f}]. \]
The cluster algebra is generated by all cluster variables in the mutation
class, whereas the upper cluster algebra is
\[ \cU(\Sigma)=\bigcap_{\Sigma'\sim\Sigma}\cL(\mathbf{x}'). \]
The Laurent phenomenon gives $\cA(\Sigma)\subseteq\cU(\Sigma)$; our
conventions follow \cite{FZ1,BFZ}.

\subsection{Theta functions and the middle algebra}

For a cluster seed, Gross--Hacking--Keel--Kontsevich construct theta functions
from the scattering diagram.  In initial seed coordinates they have the form
\[ \theta_g=\mathbf{x}^gF_g(\mathbf{y}), \qquad F_g\in1+\ZZ_{\ge0}\llbracket\mathbf{y}\rrbracket. \]
Let $D_f$ be the boundary divisor associated with a frozen index $f$,
and let $\vartheta_f^\vee$ be the corresponding theta function on the dual
cluster variety.  With the min-plus convention, the cone is
\[ \Xi =\left\{g\in\RR^m\ \middle|\ (\vartheta_f^\vee)^{\mathrm{trop}}(g)\ge0 \text{ for every frozen }f\right\}. \]
If
\[ \vartheta_f^\vee=\sum_{\alpha\in A_f}c_\alpha\mathbf{x}^\alpha, \qquad c_\alpha>0, \]
then
\[ (\vartheta_f^\vee)^{\mathrm{trop}}(g) =\min_{\alpha\in A_f}\langle\alpha,g\rangle, \]
so $\Xi$ is a rational polyhedral cone.  The middle algebra is the theta
subalgebra regular along the chosen boundary.  Under the hypotheses of
\cite{GHKK}, and in the valuative formulation of \cite{CMMM}, it has basis
\begin{equation}
\label{eq:middle-theta-basis}
 \{\theta_g\mid g\in\Xi\cap\ZZ^m\}.
\end{equation}
Valuative independence implies that regularity may be tested one theta
function at a time.

Fix a translation-invariant total order on $\ZZ^m$.  In a seed torus, the
seed valuation of a nonzero Laurent polynomial is the least exponent in its
support.  This lowest-term valuation has one-dimensional leaves: if two
Laurent polynomials have the same least exponent, subtracting the ratio of
their leading coefficients strictly increases the valuation.  This is the
standard lowest-term example in \cite[Example~2.2(2)]{KM}.

We shall also use the following elementary Khovanskii-basis observation.  It
is included to distinguish algebra generation from mere semigroup
generation.

\begin{lemma}
\label{lem:hilbert-theta-generation}
Suppose the finite theta functions in \eqref{eq:middle-theta-basis} are
homogeneous for a grading with finite-dimensional homogeneous pieces, and
admit a seed valuation with one-dimensional leaves and
$\initer(\theta_g)=\mathbf{x}^g$.  If $\mathcal H$ is the Hilbert basis
of $\Xi\cap\ZZ^m$, then the theta functions
$\{\theta_h\mid h\in\mathcal H\}$ generate the middle algebra.
\end{lemma}

\begin{proof}
Let $B$ be the algebra generated by the $\theta_h$, $h\in\mathcal H$.
For $g\in\Xi\cap\ZZ^m$, choose a decomposition
$g=h_1+\cdots+h_r$ with $h_i\in\mathcal H$.  After rescaling, the product
$\theta_{h_1}\cdots\theta_{h_r}$ and $\theta_g$ have the same initial
term $\mathbf{x}^g$.  Their difference is homogeneous of the same weight
and has strictly larger valuation.  The one-dimensional-leaf property
implies that only finitely many valuation levels occur in a fixed
finite-dimensional homogeneous component.  Repeating the subtraction
therefore terminates and expresses $\theta_g$ as an element of $B$.
Since the $\theta_g$ form a basis, $B$ is the middle algebra.
\end{proof}

\section{The affine quotient, logarithmic forms, and the negative Markov fiber}
\label{sec:slice-volume}

\subsection{Logarithmic charts on the affine quotient}

An algebraic torus with coordinates $z_1,\ldots,z_d$ carries the
logarithmic form
\[ \Omega_{\mathbf z}=\bigwedge_{i=1}^d d\!\log z_i =\frac{dz_1\wedge\cdots\wedge dz_d}{z_1\cdots z_d}. \]
It has neither zeros nor poles on the torus.  Cluster mutations preserve the
associated logarithmic top form up to sign.  In a geometric cluster
$\mathcal A$-realization, frozen coordinates are characterized by the
nonvanishing residues of this form along their zero divisors
\cite[Definition~11.2, Lemma~11.3, and Definition~11.5]{GS15}.  This is part
of the log Calabi--Yau interpretation of cluster varieties as spaces assembled
from algebraic tori by volume-preserving birational maps
\cite[\S1]{GHK15}.  Non-invertible frozen variables and the corresponding
boundary divisors in partial compactifications are treated in
\cite[\S2.2]{MandelCox}.

No seed on $Y_{322}$ is assumed.  We instead contract the logarithmic top
form of the double-invariant seed with the vector field of the last root
subgroup and restrict it to a local slice.  The contraction is horizontal and
invariant, and the slice identifies it with a rational top form
$\omega_{322}$ on $Y_{322}$.  Compare \cite[\S1.4]{GS15} for the analogous
construction of a quotient volume form by contraction along a freely acting
torus.  A seven-function chart $\mathbf z$ with
frozen product $P_F$ has the required logarithmic boundary exactly when
\[ P_F\bigwedge_{i=1}^7 d\!\log z_i=c\,\omega_{322} \qquad(c\in\kk^\times). \]
A nonzero constant implies that the seven functions are algebraically
independent and that $P_F=0$ has the required boundary multiplicities.  We
apply this test to two charts with the same frozen product.  Their transition
relation is a signed Markov exchange; adjoining the coefficient $\zeta$
then realizes it as a specialization of ordinary cluster mutations.  The
volume calculation is used only to find birational charts on $Y_{322}$ and
their frozen boundary.

\subsection{The last root derivation}

We use the cluster structure on the double-invariant algebra recalled in
Appendix~\ref{app:partial-ring}.  Write
$\Sigma_s=(\mathbf s^\circ,B_s)$.  Its initial extended cluster, ordered so
that the first three entries are mutable, is
\begin{equation}
\label{eq:double-initial-cluster}
 \mathbf s^\circ=(s_1,\ldots,s_8)
 =\left(
 s_{110,20}^{11},
 s_{100,10}^{01},
 s_{100,10}^{10},
 s_{110,11}^{02},
 s_{110,11}^{20},
 s_{200,11}^{11},
 s_{111,21}^{12},
 s_{111,21}^{21}
 \right).
\end{equation}
Here the superscripts record the first two highest weights, whereas the
subscript records the remaining $\GL_2$-weight.  Let $\partial$ be the locally nilpotent derivation induced by the positive root
subgroup of the last $\GL_2$.  The root strings through the initial
variables are
\begin{equation}
\label{eq:partial-on-double-seed}
 \partial(s_2)=s_3,
 \qquad
 \partial(s_4)=s_9,
 \qquad
 \partial(s_9)=2s_5,
 \qquad
 \partial(s_7)=-s_8,
\end{equation}
while $\partial$ kills $s_1,s_3,s_5,s_6,s_8$.  The middle coefficient in the
length-three string is
\begin{equation}
\label{eq:middle-root-coefficient}
 s_9=s_{110,11}^{11}
 =\frac{s_1s_6+s_2^2s_5+s_3^2s_4}{s_2s_3}.
\end{equation}
Thus
\[ \UR=\ker(\partial:\SR\longrightarrow\SR). \]
The signs in \eqref{eq:partial-on-double-seed} are fixed by the convention that the
positive root subgroup acts on the two tensor slices by
$(X,Y)\mapsto(X,Y+tX)$.

\subsection{A birational chart on the affine quotient}

On the principal open set $s_3\ne0$, put
\[ \xi=\frac{s_2}{s_3}. \]
Then $\partial(\xi)=1$.  Since $s_3=u_1$ is invariant, the local-slice
theorem gives
\[ \SR[s_3^{-1}]=\UR[u_1^{-1}][\xi]. \]
Thus $\pi^{-1}(D(u_1))\simeq D(u_1)\times\mathbb A^1$, equivariantly for
the translation action of $\mathbb G_a$ on the coordinate $\xi$
\cite[Chapter~1]{Freudenburg}.  In particular, this open subset has a
geometric quotient, whereas $Y_{322}$ denotes the global affine
categorical quotient.  We use the following seven invariant functions as
birational coordinates on its base:
\begin{equation}
\label{eq:quotient-coordinates}
 \begin{aligned}
 q_1&=u_3=s_1,&
 q_2&=u_1=s_3,&
 q_3&=u_2=s_5,\\
 q_4&=u_4=s_6,&
 q_5&=u_5=s_8,&
 q_6&=u_6=s_3s_9-2s_2s_5,\\
 &&q_7&=u_8=s_2s_8+s_3s_7.
 \end{aligned}
\end{equation}
Their weights are, respectively,
\[ (110;20;11),\ (100;10;10),\ (110;11;20),\ (200;11;11),\ (111;21;21),\ (210;21;21),\ (211;31;22). \]

\begin{proposition}
\label{prop:birational-quotient-chart}
The functions $\xi,q_1,\ldots,q_7$ form a triangular birational coordinate
system.  More precisely,
\begin{equation}
\label{eq:inverse-quotient-chart}
 \begin{gathered}
 s_1=q_1,\qquad s_2=\xi q_2,\qquad s_3=q_2,
 \qquad s_5=q_3,\qquad s_6=q_4,\qquad s_8=q_5,\\
 s_7=\frac{q_7}{q_2}-\xi q_5,
 \qquad
 s_4=\xi^2q_3+\frac{\xi q_6}{q_2}-\frac{q_1q_4}{q_2^2}.
 \end{gathered}
\end{equation}
Consequently,
\[ \kk(X_{\mathrm{dbl}})=\kk(q_1,\ldots,q_7)(\xi), \qquad \kk(Y_{322})=\kk(q_1,\ldots,q_7). \]
\end{proposition}

\begin{proof}
The first six formulas in \eqref{eq:inverse-quotient-chart} follow directly
from \eqref{eq:quotient-coordinates}.  Solving the last equation in
\eqref{eq:quotient-coordinates} gives the formula for $s_7$.  Substitution
of $s_2=\xi q_2$ into the formula for $q_6$, followed by
\eqref{eq:middle-root-coefficient}, gives the formula for $s_4$.  This
proves the first field equality.  Since $\partial(\xi)=1$ and all $q_i$ are
annihilated by $\partial$, the kernel of $\partial$ in this rational function field is
exactly $\kk(q_1,\ldots,q_7)$.
\end{proof}

The coordinate $\xi$ chooses a section of the trivialized
$\mathbb G_a$-bundle over $D(u_1)$.  It is used only for the calculation.
The differential form below descends by contraction with the vector field
$\partial$, and therefore does not depend on this choice of section.
The frozen variables of the double seed are $s_4,\ldots,s_8$.  Normalize
its logarithmic top form by
\begin{equation}
\label{eq:double-log-volume}
 \omega_s
 =s_4s_5s_6s_7s_8\bigwedge_{i=1}^8\dlog s_i
 =\frac{ds_1\wedge\cdots\wedge ds_8}{s_1s_2s_3}.
\end{equation}

\begin{proposition}
\label{prop:quotient-volume}
The contracted form $\omega_{322}=\iota_{\partial}\omega_s$ is
\begin{equation}
\label{eq:quotient-volume}
 \omega_{322}
 =-\frac{dq_1\wedge\cdots\wedge dq_7}{q_1q_2^3}.
\end{equation}
\end{proposition}

\begin{proof}
The divergence of $\partial$ with respect to
$ds_1\wedge\cdots\wedge ds_8$ is $s_3/s_2$, whereas
$\partial\log((s_1s_2s_3)^{-1})=-s_3/s_2$.  Hence
$\mathcal L_{\partial}\omega_s=0$.  It follows that
$\iota_{\partial}\omega_s$ is horizontal and invariant, so it descends to
$Y_{322}$.

A direct differentiation of \eqref{eq:quotient-coordinates} gives
\[ \det\frac{\partial(\xi,q_1,\ldots,q_7)} {\partial(s_1,\ldots,s_8)} =-\frac{s_3^2}{s_2}. \]
Since $\partial(\xi)=1$ and $\partial(q_i)=0$, the orientation convention in
\eqref{eq:double-log-volume} gives
$\omega_s=d\xi\wedge\omega_{322}$.  Substituting the
Jacobian and $s_1=q_1,s_3=q_2$ yields \eqref{eq:quotient-volume}.
\end{proof}

Let $\mathbf z=(z_1,\ldots,z_7)$ be a proposed rational chart on
$Y_{322}$, and let $F$ be its frozen subset, with frozen product
$P_F=\prod_{f\in F}f$.  Define
\begin{equation}
\label{eq:volume-test}
 \mathcal V_{\mathbf z,F}(q)
 =-q_1q_2^3P_F(q)
 \det\left(\frac{\partial\log z_i}{\partial q_j}\right)_{1\le i,j\le7}.
\end{equation}
Then
\[ P_F\bigwedge_{i=1}^7\dlog z_i =\mathcal V_{\mathbf z,F}(q)\,\omega_{322}. \]
Thus a logarithmic cluster-type chart must satisfy
\[ \mathcal V_{\mathbf z,F}(q)\in\kk^\times. \]
This single condition tests both algebraic independence and the boundary
multiplicity.  The term \emph{anticanonical boundary} is used here in the
logarithmic sense.  On the partial chart in which the mutable coordinates
remain invertible and the frozen functions are allowed to vanish,
$\bigwedge_i d\!\log z_i$ has simple poles along the reduced frozen
boundary $D_F=\sum_{f\in F}D_f$.  The identity
\[ P_F\bigwedge_{i=1}^7d\!\log z_i=c\,\omega_{322} \]
says on the common smooth locus that $K_{Y_{322}}+D_F$ is trivial.  Thus
$D_F$ is anticanonical there.  This is a statement about the boundary
selected by the chart; it does not assume a global cluster compactification
of $Y_{322}$.

Every $d\!\log z_i$ has torus weight zero.  Hence the character of the
frozen product equals that of the descended top form.  We call this common
character the \emph{anticanonical weight}:
\begin{equation}
\label{eq:anticanonical-weight}
 \kappa_{\mathrm{ac}}:=\sum_{f\in F}\wt(f)=\wt(\omega_{322})=(6,4,2;7,5;7,5).
\end{equation}

\subsection{The four-frozen logarithmic chart}

The quotient construction produces the following thirteen nonzero
highest-weight invariants.  We use the normalizations coming from
\eqref{eq:quotient-coordinates} and \eqref{eq:basic-relations} and write
\begin{equation}
\label{eq:thirteen-invariant-weights}
\begin{array}{lllll}
 u_1=u_{100,10}^{10},&u_2=u_{110,11}^{20},&u_3=u_{110,20}^{11},&u_4=u_{200,11}^{11},\\
 u_5=u_{111,21}^{21},&u_6=u_{210,21}^{21},&u_7=u_{211,22}^{31},&u_8=u_{211,31}^{22},\\
 u_9=u_{220,22}^{22},&u_{10}=u_{221,32}^{32},&u_{11}=u_{321,33}^{42},&u_{12}=u_{321,42}^{33},&u_{13}=u_{222,33}^{33}.
\end{array}
\end{equation}
The normalizations of $u_1,\ldots,u_8$ agree with
\eqref{eq:quotient-coordinates}; the remaining normalizations are fixed by
the relations in \eqref{eq:basic-relations} below.  The multiplicity-one
statements used later are proved, after generation is established, in
\cref{lem:relevant-one-dimensional-weights}.

Among the thirteen weights in \eqref{eq:thirteen-invariant-weights}, the
anticanonical weight $\kappa_{\mathrm{ac}}$ in
\eqref{eq:anticanonical-weight} has precisely three
square-free decompositions:
\[ \{u_{11},u_{12}\}, \qquad \{u_7,u_8,u_9\}, \qquad \{u_2,u_3,u_4,u_{13}\}. \]
The last one is the four-factor branch used below; it supports a full-rank logarithmic chart with three mutable variables.

\begin{proposition}
\label{prop:volume-selected-seed}
The two charts on $Y_{322}$
\[ \Sigma_4=(u_1,u_5,u_7\,;u_2,u_3,u_4,u_{13}), \qquad \Sigma_4'=(u_1,u_5,u_8\,;u_2,u_3,u_4,u_{13}) \]
contain the four-factor anticanonical boundary
$F_4=\{u_2,u_3,u_4,u_{13}\}$ and satisfy, in the displayed variable
order,
\begin{equation}
\label{eq:two-volume-identities}
 \begin{aligned}
 u_2u_3u_4u_{13}\bigwedge_{z\in\Sigma_4}\dlog z
 &=-\omega_{322},\\
 u_2u_3u_4u_{13}\bigwedge_{z\in\Sigma_4'}\dlog z
 &=\omega_{322}.
 \end{aligned}
\end{equation}
In particular, both are algebraically independent logarithmic charts with
the same anticanonical boundary.
\end{proposition}

\begin{proof}
Using \eqref{eq:quotient-coordinates}, one has
\[ u_7=\frac{q_5q_6+q_3q_7}{q_1}, \qquad u_{13}=\frac{q_3q_7^2+q_5q_6q_7-q_1q_4q_5^2}{q_1q_2^2}. \]
Substitution in \eqref{eq:volume-test} gives respectively the constants
$-1$ and $1$, proving \eqref{eq:two-volume-identities}.  These Jacobian identities prove the proposition.
\end{proof}

The opposite constants are the Jacobian signs for adjacent logarithmic
charts.  Weight homogeneity then leaves two monomials in the
exchange relation, and the fixed normalizations give
\[ u_7u_8=u_4u_5^2+u_1^2u_{13}. \]
The remaining two one-step exchanges are
\begin{equation}
\label{eq:first-two-signed-exchanges}
 \begin{aligned}
 u_1u_1'&=u_2u_4u_5^2+u_3u_7^2,\\
 u_5u_5'&=u_1^2u_2u_{13}-u_3u_7^2.
 \end{aligned}
\end{equation}
The contracted logarithmic form verifies the adjacent charts, and the grading
determines the signed exchanges.

The sign in the second exchange cannot be removed simultaneously with the
other two signs by rescaling the seven functions.  We therefore place the
three exchanges in an ordinary coefficient family rather than introduce a
separate signed mutation theory.  Put
\[ K_\zeta=\kk[u_2,u_3,u_4,u_{13},\zeta^{\pm1}] \]
and let $\Sigma_{u,\zeta}$ have mutable cluster $(u_1,u_5,u_7)$
and frozen coefficients $(u_2,u_3,u_4,u_{13},\zeta)$.  We retain the
seven-variable order on the negative fiber,
\begin{equation}
\label{eq:markov-variable-order}
 \mathbf u^\circ=(u_1,u_5,u_7\,;u_2,u_3,u_4,u_{13}),
\end{equation}
and order the variables of the ordinary family by
\begin{equation}
\label{eq:zeta-variable-order}
 \mathbf u_\zeta^\circ
 =(u_1,u_5,u_7\,;u_2,u_3,u_4,u_{13},\zeta).
\end{equation}
Its extended exchange matrix, in the preferred
$\mu\times(\mu+\mathrm{frozen})$ convention, is
\begin{equation*}
 B_{u,\zeta}=
 \begin{bmatrix}
  0&-2& 2&-1& 1&-1& 0&0\\
  2& 0&-2& 1&-1& 0& 1&1\\
 -2& 2& 0& 0& 0& 1&-1&0
 \end{bmatrix}.
\end{equation*}
The rows are indexed by $u_1,u_5,u_7$, and the columns are ordered by
\eqref{eq:zeta-variable-order}.  Thus $\Sigma_{u,\zeta}$ is an ordinary
cluster seed over $K_\zeta$ with exchanges
\begin{equation}
\label{eq:zeta-family-exchanges}
 \begin{aligned}
 u_1u_{1,\zeta}'&=u_2u_4u_5^2+u_3u_7^2,\\
 u_5u_{5,\zeta}'&=u_3u_7^2+\zeta u_1^2u_2u_{13},\\
 u_7u_{7,\zeta}'&=u_4u_5^2+u_1^2u_{13}.
 \end{aligned}
\end{equation}
On the coefficient fiber $\zeta=-1$, the first and third exchanges are
unchanged and the second becomes the negative of the signed exchange in
\eqref{eq:first-two-signed-exchanges}.  Multiplication of the adjacent
variable by the unit $-1$ does not change its Laurent ring.  We call this
specialization the \emph{negative Markov fiber}.


The orthogonal $D_5$-parabolic origin of this four-frozen seed, including
an integral Gale-dual interpretation of its exchange matrix, is explained in
Appendix~\ref{app:D5-parabolic}.

\section{The invariant, upper, and middle algebras}
\label{sec:ring-equalities}

\subsection{The invariant ring and the four-chart intersection}
\label{subsec:invariant-upper}

We first use two slices to prove $\UR=\Ugen$.  We then show
$\Uneg\subseteq\UR$ and $\Mu\subseteq\Uneg$.  The theta Hilbert
basis places twelve invariant generators in $\Mu$, and two discriminant
identities place the remaining generator in $\Uneg$.  Together these
steps give
\[ \UR=\Uneg=\Mu[\udisc], \]
and reduce each fixed triple-weight space to one theta level.

Put
\[ \Ugen=\kk[u_1,\ldots,u_{13}]\subseteq\UR. \]
The following six identities will be used throughout:
\begin{equation}
\label{eq:basic-relations}
 \begin{alignedat}{2}
 u_3u_7&=u_2u_8+u_5u_6,
 &\qquad u_1^2u_9&=u_6^2+4u_2u_3u_4,\\
 u_1u_{10}&=2u_2u_8+u_5u_6,
 &u_1u_{11}&=-2u_2u_4u_5-u_6u_7,\\
 u_1u_{12}&=2u_3u_4u_5-u_6u_8,
 &u_1^2u_{13}&=u_7u_8-u_4u_5^2.
 \end{alignedat}
\end{equation}
They may be checked in the quotient coordinates of
\cref{prop:birational-quotient-chart}.  We prove that no further invariant
generators are needed.

Choose the generating set $s_1,\ldots,s_{14}$ of $\SR$ recorded in
\cref{tab:partial-generators}; it is obtained from the minimal generating
set in Appendix~\ref{app:partial-ring} by adjoining one redundant one-step
cluster variable.  The nonzero values of $\partial$ are
\begin{equation}
\label{eq:partial-on-s-generators}
 \begin{gathered}
 \partial(s_2)=s_3,\qquad \partial(s_4)=s_9,\qquad
 \partial(s_7)=-s_8,\qquad \partial(s_9)=2s_5,\\
 \partial(s_{10})=s_{11},\qquad
 \partial(s_{12})=2s_{10},\qquad
 \partial(s_{13})=2s_5s_2.
 \end{gathered}
\end{equation}
The remaining seven generators are killed by $\partial$.  In terms of
this set, the thirteen invariants are
\begin{equation}
\label{eq:u-in-terms-of-s}
 \begin{array}{lll}
 u_1=s_3,&u_2=s_5,&u_3=s_1,\\
 u_4=s_6,&u_5=s_8,&u_7=s_{11},\\
 u_{13}=s_{14},&u_6=s_3s_9-2s_5s_2,&u_8=s_3s_7+s_2s_8,\\
 u_9=s_9^2-4s_5s_4,&u_{10}=s_9s_8+2s_5s_7,&u_{11}=2s_5s_{10}-s_9s_{11},\\
 \multicolumn{3}{l}{u_{12}=s_2s_9s_8-s_3s_9s_7+2s_2s_5s_7-2s_3s_4s_8.}
 \end{array}
\end{equation}
The root strings \eqref{eq:partial-on-s-generators} show directly that all
expressions in \eqref{eq:u-in-terms-of-s} lie in $\ker\partial$.

We use the two slices
\[ \xi_1=\frac{s_2}{s_3}=\frac{s_2}{u_1}, \qquad \xi_2=\frac{s_9}{2s_5}=\frac{s_9}{2u_2}. \]
Both satisfy $\partial(\xi_i)=1$.  They come from the shortest linear
and quadratic root strings, and their denominators are the two lowest-degree
nonassociate highest coefficients.  For $i=1,2$, set
\[ \rho_i(g)=\exp(-\xi_i\partial)g =\sum_{j\ge0}\frac{(-\xi_i)^j}{j!}\partial^j(g). \]
The sum is finite because $\partial$ is locally nilpotent.  The
nontrivial values of the two invariantization maps are
\begin{equation}
\label{eq:invariantization-table}
\begin{array}{c|c|c}
 g&\rho_1(g)&\rho_2(g)\\ \hline
 s_{10}&-u_4u_5/u_1&u_{11}/(2u_2)\\
 s_7&u_8/u_1&u_{10}/(2u_2)\\
 s_{12}&u_4u_8/u_1^2&u_7u_9/(4u_2^2)\\
 s_{13}&u_3u_4/u_1&u_1u_9/(4u_2)\\
 s_2&0&-u_6/(2u_2)\\
 s_4&-u_3u_4/u_1^2&-u_9/(4u_2)\\
 s_9&u_6/u_1&0
\end{array}
\end{equation}
The generators $s_1,s_3,s_5,s_6,s_8,s_{11},s_{14}$ are already
invariant.  The local-slice theorem \cite[Chapter~1]{Freudenburg} and \eqref{eq:invariantization-table}
therefore give
\begin{equation}
\label{eq:two-localized-equalities}
 \UR[u_1^{-1}]=\Ugen[u_1^{-1}],
 \qquad
 \UR[u_2^{-1}]=\Ugen[u_2^{-1}].
\end{equation}

\begin{lemma}[Two-localization lemma]
\label{lem:two-localization}
Let $A$ be a domain and let $f,g\in A$ be nonzero.  If multiplication
by $g$ is injective on $A/fA$, then, inside $\Frac(A)$,
\[ A[f^{-1}]\cap A[g^{-1}]=A. \]
No Noetherian or normality hypothesis is required.
\end{lemma}

\begin{proof}
Write an element of the intersection as
$x=a/f^p=b/g^q$.  Then $g^qa=f^pb$.  Modulo $f$, injectivity of
multiplication by $g^q$ gives $a\in fA$.  Cancelling one power of
$f$ and iterating removes the entire denominator $f^p$, so $x\in A$.
\end{proof}

\begin{lemma}
\label{lem:Ugen-regular-sequence}
The pair $u_1,u_2$ is a regular sequence in $\Ugen$.  Equivalently,
multiplication by $u_2$ is injective on $\Ugen/u_1\Ugen$.
\end{lemma}

\begin{proof}
The algebra $\Ugen$ is a domain, so $u_1$ is a non-zero-divisor.  The
single colon computation in Appendix~\ref{app:regular-sequence} gives
\[ (I+(U_1)):U_2=I+(U_1) \]
for the defining ideal $I$ of $\Ugen$, which is the asserted
injectivity modulo $u_1$.
\end{proof}

\begin{theorem}
\label{thm:thirteen-generate}
The thirteen functions in \eqref{eq:thirteen-invariant-weights} generate the
triple-invariant algebra:
\[ \UR=\Ugen. \]
\end{theorem}

\begin{proof}
By \cref{lem:two-localization,lem:Ugen-regular-sequence},
\begin{equation}
\label{eq:Ugen-localization-intersection}
 \Ugen[u_1^{-1}]\cap\Ugen[u_2^{-1}]=\Ugen.
\end{equation}
Let $f\in\UR$.  The two slice equalities
\eqref{eq:two-localized-equalities} place $f$ in both rings on the
left-hand side of \eqref{eq:Ugen-localization-intersection}; hence
$f\in\Ugen$.  The reverse inclusion is tautological.
\end{proof}

\begin{lemma}
\label{lem:relevant-one-dimensional-weights}
For each weight in the following table, the monomial is the unique
monomial in $u_1,\ldots,u_{13}$ of that weight:
\begin{equation}
\label{eq:unique-weight-monomials}
\begin{array}{c|c@{\qquad}c|c}
\text{weight}&\text{unique monomial}&\text{weight}&\text{unique monomial}\\ \hline
(210;21;21)&u_6&(221;32;32)&u_{10}\\
(321;33;42)&u_{11}&(321;42;33)&u_{12}\\
(320;32;32)&u_9u_1&(331;43;43)&u_9u_5\\
(431;44;53)&u_9u_7&(431;53;44)&u_9u_8
\end{array}
\end{equation}
Consequently, the corresponding homogeneous components of $\UR$ are
one-dimensional.
\end{lemma}

\begin{proof}
By \cref{thm:thirteen-generate}, every homogeneous element of $\UR$ is a
linear combination of monomials in the thirteen generators.  For each row
of \eqref{eq:unique-weight-monomials}, solving the seven nonnegative integer
equations obtained from the weights in
\eqref{eq:thirteen-invariant-weights} gives the exponent
vector.  Thus the homogeneous component is spanned by that nonzero
monomial.  The \texttt{weights} task in \path{Kron322.py} checks
these eight semigroup searches.
\end{proof}

We compare $\UR$ with the first four cluster charts of the ordinary
coefficient family.  Let
\[ K_u=\kk[u_2,u_3,u_4,u_{13}]. \]
Specializing the one-step variables of \eqref{eq:zeta-family-exchanges}, set
\[ u_1'=u_{1,\zeta}'\big|_{\zeta=-1},\qquad u_5'=-u_{5,\zeta}'\big|_{\zeta=-1},\qquad u_7'=u_{7,\zeta}'\big|_{\zeta=-1}. \]
They are regular invariants; the basic relations give
\begin{equation}
\label{eq:partners-in-u}
 u_1'=u_7u_{10}-u_1u_2u_{13},
 \qquad
 u_5'=u_1u_{11}+u_2u_4u_5,
 \qquad
 u_7'=u_8.
\end{equation}

Let $\cL_\circ(\zeta),\cL_1(\zeta),\cL_5(\zeta),\cL_7(\zeta)$
be the Laurent rings of the initial seed $\Sigma_{u,\zeta}$ and its
three one-step mutations:
\begin{equation*}
 \begin{aligned}
 \cL_\circ(\zeta)&=K_\zeta[u_1^{\pm1},u_5^{\pm1},u_7^{\pm1}],\\
 \cL_1(\zeta)&=K_\zeta[(u_{1,\zeta}')^{\pm1},u_5^{\pm1},u_7^{\pm1}],\\
 \cL_5(\zeta)&=K_\zeta[u_1^{\pm1},(u_{5,\zeta}')^{\pm1},u_7^{\pm1}],\\
 \cL_7(\zeta)&=K_\zeta[u_1^{\pm1},u_5^{\pm1},(u_{7,\zeta}')^{\pm1}].
 \end{aligned}
\end{equation*}
Write
\[ \cL_t^-:=\cL_t(\zeta)/(\zeta+1)\cL_t(\zeta) \]
for the specialization at $\zeta=-1$.  Since $-1$ is a unit,
these four rings are identified with
\begin{equation*}
 \begin{aligned}
 \cL_\circ^-&=K_u[u_1^{\pm1},u_5^{\pm1},u_7^{\pm1}],\\
 \cL_1^-&=K_u[(u_1')^{\pm1},u_5^{\pm1},u_7^{\pm1}],\\
 \cL_5^-&=K_u[u_1^{\pm1},(u_5')^{\pm1},u_7^{\pm1}],\\
 \cL_7^-&=K_u[u_1^{\pm1},u_5^{\pm1},(u_7')^{\pm1}].
 \end{aligned}
\end{equation*}
The \texttt{laurent} task in \path{Kron322.py} writes
explicit Laurent expressions for all thirteen $u_i$ in these four charts
and verifies the six relations in \eqref{eq:basic-relations}.
We define the \emph{four-chart intersection on the negative fiber} by
\[ \Uneg:=\cL_\circ^-\cap\cL_1^-\cap\cL_5^-\cap\cL_7^- \subseteq\Frac(\UR). \]
These four charts are the initial chart and its three codimension-one
neighbors, so their intersection is the finite upper bound determined by the chosen
logarithmic seed.  The definition uses only ordinary cluster
charts before specialization.  We do not assert that specialization commutes
with the infinite intersection defining the full upper cluster algebra
$\cU(\Sigma_{u,\zeta})$; no such base-change statement is needed below.

\begin{lemma}
\label{lem:four-chart-coprimality}
The elements $u_1,u_5,u_7$ are pairwise nonassociate primes in the UFD
$\UR$, and
\[ \gcd(u_1,u_1')=\gcd(u_5,u_5')=\gcd(u_7,u_7')=1. \]
\end{lemma}

\begin{proof}
The algebra $\UR$ is factorial by the factoriality theorem for invariants
of a factorial affine variety under a connected group with trivial character
group, applied to $U_3\times U_2\times U_2$; see \cite{PV}.  Its only
units are scalars, by the positive polynomial-degree grading.

Because the grading torus is connected, every irreducible factor of a
homogeneous semi-invariant in the UFD $\UR$ is again a homogeneous
semi-invariant.  The only semigroup weight of polynomial degree one is
$\wt(u_1)$.  Hence a factorization of $u_5$ would have degree split
$1+2$, but
\[ \wt(u_5)-\wt(u_1)=(011;11;11) \]
is not dominant.  A factorization of $u_7$ has split $1+3$ or
$2+2$.  The first is excluded by
\[ \wt(u_7)-\wt(u_1)=(111;12;21). \]
The degree-two semigroup weights and the corresponding differences are
\[
\begin{array}{c|c}
\text{degree-two weight }\eta&\wt(u_7)-\eta\\ \hline
2\wt(u_1)&(011;02;11)\\
\wt(u_2)&(101;11;11)\\
\wt(u_3)&(101;02;20)\\
\wt(u_4)&(011;11;20).
\end{array}
\]
Every difference is nondominant.  Thus $u_1,u_5,u_7$ are irreducible,
hence prime, and their distinct weights make them pairwise nonassociate.

Modulo these primes, \eqref{eq:partners-in-u} gives
\[ u_1'\equiv u_7u_{10}\pmod{u_1},\qquad u_5'\equiv u_1u_{11}\pmod{u_5},\qquad u_7'=u_8. \]
The relevant weight differences are
\[
\begin{array}{c|c}
\text{putative quotient}&\text{weight difference}\\ \hline
u_7/u_1&(111;12;21)\\
u_{10}/u_1&(121;22;22)\\
u_{11}/u_5&(210;12;21)\\
u_8/u_7&(000;1,-1;-1,1).
\end{array}
\]
None belongs to the dominant weight semigroup.  Therefore the indicated
prime divides neither factor on the right, proving the three gcd statements.
\end{proof}

\begin{proposition}
\label{prop:upper-contained-in-invariants}
One has $\Uneg\subseteq\UR$.
\end{proposition}

\begin{proof}
Let $f\in\Uneg$.  From the initial chart,
\[ f\in\UR[u_1^{-1},u_5^{-1},u_7^{-1}]. \]
In a reduced UFD expression for $f$, the only possible irreducible
denominator factors are therefore $u_1,u_5,u_7$.  On the other hand,
\[ f\in\UR[(u_1')^{-1},u_5^{-1},u_7^{-1}]. \]
By \cref{lem:four-chart-coprimality}, all three inverted elements in this
second localization are coprime to $u_1$; hence the $u_1$-adic exponent
of $f$ is nonnegative.  The $u_5'$- and $u_7'$-charts similarly show
that the $u_5$- and $u_7$-adic exponents are nonnegative.  Thus the
reduced denominator has no irreducible factor, and $f\in\UR$.
\end{proof}

\subsection{The middle algebra and the discriminant}
\label{subsec:middle-extension}

We distinguish three operations: passage from principal to geometric
coefficients, specialization at $\zeta=-1$, and extension across the four
noninvertible coefficient divisors.  Put
\[ P=\kk[y_1^{\pm1},y_2^{\pm1},y_3^{\pm1}], \qquad K_\zeta^{\mathrm{tor}}=\kk[u_2^{\pm1},u_3^{\pm1},u_4^{\pm1},u_{13}^{\pm1},\zeta^{\pm1}] \]
and
\[ K_u^{\mathrm{tor}} =\kk[u_2^{\pm1},u_3^{\pm1},u_4^{\pm1},u_{13}^{\pm1}]. \]
Let $\Sigma_{\mathrm{Mar}}^{\mathrm{prin}}$ be the principal-coefficient
seed with the Markov principal part.  The relevant coefficient maps are
\[ P\xrightarrow{\ \varphi_\zeta\ }K_\zeta^{\mathrm{tor}} \xrightarrow{\ \operatorname{ev}_{-1}\ }K_u^{\mathrm{tor}}, \qquad \operatorname{ev}_{-1}(\zeta)=-1, \]
where, in normalized $y$-coordinates,
\begin{equation}
\label{eq:principal-coefficient-specialization}
 \varphi_\zeta(y_1)=\frac{u_3}{u_2u_4},\qquad
 \varphi_\zeta(y_2)=\frac{\zeta u_2u_{13}}{u_3},\qquad
 \varphi_\zeta(y_3)=\frac{u_4}{u_{13}}.
\end{equation}
Clearing the common frozen units in the three exchanges gives
\eqref{eq:zeta-family-exchanges}.

Let
\[ N_{\mathrm{mut}}=\ZZ e_1\oplus\ZZ e_2\oplus\ZZ e_3, \qquad L_\partial=\bigoplus_{i=4}^7\ZZ e_i, \qquad L_{\mathrm{fr}}=L_\partial\oplus\ZZ e_\zeta. \]
In the ordered frozen basis
$(e_4,e_5,e_6,e_7,e_\zeta)$, corresponding to
$(u_2,u_3,u_4,u_{13},\zeta)$, the monomial map
$\varphi_\zeta$ is induced by
\begin{equation}
\label{eq:coefficient-exponent-map}
 \begin{aligned}
 \kappa:\ZZ^3&\longrightarrow L_{\mathrm{fr}},\\
 \epsilon_1&\longmapsto(-1,1,-1,0,0),\\
 \epsilon_2&\longmapsto(1,-1,0,1,1),\\
 \epsilon_3&\longmapsto(0,0,1,-1,0).
 \end{aligned}
\end{equation}
Together with the identity on mutable directions, this gives the diagram of
label lattices
\begin{equation}
\label{eq:coefficient-lattice-diagram}
 \begin{array}{ccc}
 N_{\mathrm{mut}}\oplus\ZZ^3
 &\xrightarrow{\ \mathrm{id}\oplus\kappa\ }&
 N_{\mathrm{mut}}\oplus L_{\mathrm{fr}}\\[-1mm]
 \big\downarrow&&\big\downarrow\\[-1mm]
 N_{\mathrm{mut}}&\xrightarrow{\ \mathrm{id}\ }&N_{\mathrm{mut}},
 \end{array}
\end{equation}
where the vertical arrows forget coefficient directions.  Put
\[
 N_{u,\zeta}=N_{\mathrm{mut}}\oplus L_{\mathrm{fr}},
 \qquad
 N_u=N_{\mathrm{mut}}\oplus L_\partial.
\]
For $g\in N_u$, let $\widetilde\theta_{(g,0)}$ be the geometric-coefficient
theta function with zero $e_\zeta$-coordinate, and put
\[ \theta_g^-:=\operatorname{ev}_{-1} \bigl(\widetilde\theta_{(g,0)}\bigr). \]

\begin{lemma}[Principal-to-geometric theta functions]
\label{lem:principal-to-geometric-theta}
After the frozen rescaling used to pass from normalized $y$-coordinates to
\eqref{eq:zeta-family-exchanges}, the map
$\mathrm{id}\oplus\kappa$ in
\eqref{eq:coefficient-lattice-diagram} is a linear morphism of seed data.
If a principal-coefficient theta function and its geometric-coefficient image
are normalized to have the same initial mutable monomial, then
\begin{equation}
\label{eq:principal-geometric-theta-map}
 \varphi_\zeta\bigl(\vartheta^{\mathrm{prin}}_{(a,0)}\bigr)
 =\widetilde\theta_{(a,0)}
 \qquad(a\in N_{\mathrm{mut}}).
\end{equation}
Moreover, for $q\in L_{\mathrm{fr}}$, frozen translation gives
\begin{equation}
\label{eq:frozen-theta-translation}
 \widetilde\theta_{(a,q)}=X^q\widetilde\theta_{(a,0)}.
\end{equation}
Consequently, finiteness of all principal-coefficient theta functions implies
finiteness of all theta functions of $\Sigma_{u,\zeta}$ on the coefficient
torus.
\end{lemma}

\begin{proof}
The exponent vectors in \eqref{eq:coefficient-exponent-map} are those
of the three monomials in
\eqref{eq:principal-coefficient-specialization}.  In the mutable-row
convention, the principal extended matrix is
$[B_{\mathrm{Mar}}\mid I_3]$, whereas
\[ B_{u,\zeta}=[B_{\mathrm{Mar}}\mid\kappa^{\mathsf T}]. \]
Thus, after clearing the common frozen units,
$\mathrm{id}\oplus\kappa$ satisfies the defining seed-datum identities
and is the linear morphism inducing $\varphi_\zeta$.  Functoriality of
theta functions under linear morphisms of seed data
\cite[Theorem~4.9]{CMMM} gives
\eqref{eq:principal-geometric-theta-map}; the stated normalization removes
the frozen monomial introduced by the rescaling.  A frozen character has no
walls, so multiplication by $X^q$ translates the theta index and gives
\eqref{eq:frozen-theta-translation}.  The finiteness statement follows.
\end{proof}

\begin{proposition}[Specialization on the coefficient torus]
\label{prop:specialized-theta-span}
Every $\widetilde\theta_{(g,r)}$ is a finite Laurent polynomial.  The
functions $\theta_g^-$, $g\in N_u$, are valuatively and linearly
independent, and their $\kk$-span is an algebra.  If
\[ \cL_t^{-,\mathrm{tor}} :=\cL_t^-[(u_2u_3u_4u_{13})^{-1}], \qquad t=\circ,1,5,7, \]
then
\[ \theta_g^-\in \bigcap_{t=\circ,1,5,7}\cL_t^{-,\mathrm{tor}} \qquad(g\in N_u). \]
\end{proposition}

\begin{proof}
For the once-punctured-torus Markov seed with principal coefficients, Zhou
proves the full Fock--Goncharov conjecture
\cite[Theorem~1.1]{Zhou}.  Thus every integral mutable label gives a finite
theta Laurent polynomial, the theta parametrization uses the full integral
tropical lattice, and the middle and canonical theta algebras coincide.  By
\cref{lem:principal-to-geometric-theta}, every theta function of
$\Sigma_{u,\zeta}$ on the coefficient torus is therefore finite.

Apply $\operatorname{ev}_{-1}$.  Every coefficient monomial in
$K_\zeta^{\mathrm{tor}}$ maps to a nonzero unit of
$K_u^{\mathrm{tor}}$, hence to a non-zero-divisor.  Cluster scattering
functions satisfy Assumption~3.9 of \cite{CMMM}, so
\cite[Theorem~3.10]{CMMM} shows that the finite theta functions retain their
mutable tropicalizations and remain valuatively and linearly independent.
Writing a label in $N_{u,\zeta}$ as $(g,r)$, with
$g\in N_u$ and $r\in\ZZ$ denoting the coefficient of $e_\zeta$,
frozen translation gives
\[ \operatorname{ev}_{-1}\bigl(\widetilde\theta_{(g,r)}\bigr) =(-1)^r\theta_g^-. \]
Thus the image of the ordinary middle cluster algebra, which here equals the
canonical cluster algebra, is the span of the $\theta_g^-$ on the coefficient
torus and is an algebra.  Finally, every finite theta function is Laurent in every ordinary
cluster chart; specialization gives the asserted membership in the four
coefficient-torus Laurent rings.
\end{proof}

Let $D_4,D_5,D_6,D_7$ be the four frozen boundary divisors, and define
\[ \Mu :=\operatorname{span}_{\kk}\left\{ \theta_g^-\ \middle|\ \ord_{D_f}(\theta_g^-)\ge0\text{ for }4\le f\le7 \right\}. \]

\begin{proposition}
\label{prop:negative-fiber-theta-algebra}
The algebra $\Mu$ has basis consisting of the individually regular
$\theta_g^-$, and
\[ \Mu\subseteq\Uneg. \]
\end{proposition}

\begin{proof}
By \cref{prop:specialized-theta-span}, the span of all specialized theta
functions is an algebra and the functions are valuatively independent.
Hence the order of a finite linear combination along any $D_f$ is the
minimum of the orders of its nonzero theta terms.  Such a combination is
regular along all four boundaries exactly when each theta term is regular.
This proves the basis assertion; closure under multiplication follows
because products of regular functions are regular.

On each cluster chart, $\ord_{D_f}$ is the minimum exponent of the frozen
variable $u_f$.  A theta function regular along all four $D_f$ has no
negative frozen exponent in any of the four Laurent charts.  The last
assertion of \cref{prop:specialized-theta-span} therefore improves from the
coefficient-torus rings to the nonlocalized rings $\cL_t^-$, proving
$\Mu\subseteq\Uneg$.
\end{proof}

From now on we suppress the superscript and write $\theta_g=\theta_g^-$.
Let $e_1,\ldots,e_7$ be the standard basis in the variable order
\eqref{eq:markov-variable-order}.  The four mirror boundary theta functions
have the following exponent supports:
\begin{equation*}
 \begin{aligned}
 A_4={}&\{e_4,e_1+e_4,e_1+e_3+e_4,
          e_1+2e_3+e_4,e_1+e_2+2e_3+e_4\},\\
 A_5={}&\{e_5,e_2+e_5,e_1+e_2+e_5\},\\
 A_6={}&\{e_6,e_1+e_6,e_1+e_3+e_6\},\\
 A_7={}&\{e_7,e_3+e_7,e_2+e_3+e_7\}.
 \end{aligned}
\end{equation*}
Consequently the cone is
\begin{equation}
\label{eq:explicit-cover-cone}
 \Xi=\left\{g\in\RR^7\ \middle|\
 \langle\alpha,g\rangle\ge0
 \text{ for every }\alpha\in A_4\cup A_5\cup A_6\cup A_7\right\}.
\end{equation}

\begin{lemma}\label{lem:boundary-support-computation}
After specialization at $\zeta=-1$, the four boundary theta functions tropicalize to
\[ (\vartheta_f^\vee)^{\mathrm{trop}}(g) =\min_{\alpha\in A_f}\langle\alpha,g\rangle, \qquad 4\le f\le7. \]
Consequently, the specialized middle algebra has theta basis
\[ \{\theta_g\mid g\in\Xi\cap\ZZ^7\}. \]
\end{lemma}

\begin{proof}
Work first in the ordinary family.  Mutation equivariance of broken lines
\cite[Proposition~3.6]{GHKK} transports a straight boundary broken line from
an adjacent chamber to the initial chamber.  The required reverse chamber
words and boundary-column entries are
\[
\begin{array}{c|c|c|l}
 f&\text{word}&\text{successive }b_{k,f}
 &X_f^{-1}\vartheta_f^\vee\\ \hline
 4&(1,3,2)&(-1,-2,1)
  &1+X_1+2X_1X_3+X_1X_3^2+X_1X_2X_3^2\\
 5&(2,1)&(-1,-1)&1+X_2+X_1X_2\\
 6&(1,3)&(-1,-1)&1+X_1+X_1X_3\\
 7&(3,2)&(-1,-1)&1+X_3+X_2X_3.
\end{array}
\]
Indeed, after projecting to the mutable exponent lattice, a bend of
multiplicity $r$ in direction $k$ replaces an exponent $\alpha$ by
the support of $X^\alpha(1+X_k)^r$.  The only bend of multiplicity two is
the middle step in the first row.  Thus the table gives the mutable supports
$A_f-e_f$ before coefficient specialization.

Exact monomial support need not survive a coefficient specialization.  What
is needed here is only the mutable tropicalization, and
\cite[Theorem~3.10]{CMMM} shows that it is unchanged under the
composite specialization used above.  Restoring the fixed boundary monomial
$X_f$ gives
\[ (\vartheta_f^\vee)^{\mathrm{trop}}(g) =\min_{\alpha\in A_f}\langle\alpha,g\rangle. \]
Theta reciprocity \cite[Theorem~5.19]{CMMM} identifies this tropicalization
with the boundary valuation:
\[ \ord_{D_f}(\theta_g) =\min_{\alpha\in A_f}\langle\alpha,g\rangle. \]
Thus $\theta_g$ is regular along every frozen boundary exactly when
$g\in\Xi$.  The basis assertion now follows from
\cref{prop:negative-fiber-theta-algebra}.
\end{proof}

\begin{lemma}
\label{lem:Xi-Hilbert-basis}
The Hilbert basis of $\Xi\cap\ZZ^7$ consists of the following sixteen
vectors:
\begin{equation}
\label{eq:Xi-Hilbert-basis}
 \begin{aligned}
 \mathcal H_1={}&\{e_1,e_2,e_3,2e_1-e_3+e_7\},\\
 \mathcal H_0={}&\{e_4,e_5,e_6,e_7,
 -e_1+e_2+e_4+e_6,
 -e_2+e_3+e_5, e_1-e_2+e_5+e_7,
 e_1-e_3+e_4+e_7\},\\
 \mathcal H_{-1}={}&\{-e_1+e_4+e_5+e_6,
 -e_2+e_4+e_5+e_7, -e_3+2e_4+e_6+e_7,
 -2e_2+e_3+2e_5+e_7\}.
 \end{aligned}
\end{equation}
The subscript is the mutable degree $g_1+g_2+g_3$.
\end{lemma}

\begin{proof}
Write $g=(m_1,m_2,m_3,d_4,d_5,d_6,d_7)$.  The inequalities defining
$\Xi$ are equivalent to
\[
 \begin{aligned}
 d_4&\ge\max\{0,-m_1,-m_1-m_3,-m_1-2m_3,
                    -m_1-m_2-2m_3\},\\
 d_5&\ge\max\{0,-m_2,-m_1-m_2\},\\
 d_6&\ge\max\{0,-m_1,-m_1-m_3\},\\
 d_7&\ge\max\{0,-m_3,-m_2-m_3\}.
 \end{aligned}
\]
Let $\psi\colon\RR^3\to\RR^4$ be the vector of the four maximum
functions.  Then $\Xi$ is the integral epigraph of $\psi$: every point
is a minimal lift $(m,\psi(m))$, plus a nonnegative combination of the
vertical rays $e_4,\ldots,e_7$.

A comparison of the affine forms shows that the common fan of linearity of
$\psi$ is obtained from the orthant fan by the successive star
subdivisions along
\[
 \begin{gathered}
 \langle-e_1,e_2\rangle,\qquad
 \langle e_1,-e_3\rangle,\qquad
 \langle e_1,e_1-e_3\rangle,\\
 \langle-e_2,e_3\rangle,\qquad
 \langle e_1,-e_2\rangle,\qquad
 \langle-e_2,-e_2+e_3\rangle.
 \end{gathered}
\]
The order is relevant: each cone is present when it is subdivided.
Every new ray is the sum of the primitive generators of a unimodular
two-dimensional cone, so unimodularity is preserved.  The result is a complete
fan with twenty maximal cones on which $\psi$ is linear.  The \texttt{fan}
task in \path{Kron322.py} lists the cones and checks their
determinants, active linear forms, and twelve nonvertical lifts.

On a unimodular cone every lattice point is a nonnegative integral
combination of its primitive rays, and $\psi$ is linear.  Taking minimal
lifts therefore expresses every lattice point of $\Xi$ as a nonnegative
integral combination of the sixteen vectors.  A minimal lift of a
primitive fan ray is indecomposable: equality in the subadditivity
inequalities would force both projected summands into one cone of linearity
and hence onto the same extremal ray; primitivity then forces one summand to
have zero projection, and minimality forces it to vanish.  The four vertical
unit rays are indecomposable.  Hence the sixteen vectors form the
Hilbert basis.
\end{proof}

To pass from theta indices to triple weights, order the rows by
$(u_1,u_5,u_7;u_2,u_3,u_4,u_{13})$ and the columns by
$(\lambda_1,\lambda_2,\lambda_3;\mu_1,\mu_2;\nu_1,\nu_2)$.  The
resulting weight matrix is
\begin{equation}
\label{eq:weight-matrix-W}
 W=
 \left[\begin{array}{ccc|cc|cc}
 1&0&0&1&0&1&0\\
 1&1&1&2&1&2&1\\
 2&1&1&2&2&3&1\\
 1&1&0&1&1&2&0\\
 1&1&0&2&0&1&1\\
 2&0&0&1&1&1&1\\
 2&2&2&3&3&3&3
 \end{array}\right].
\end{equation}
Thus $\wt(\theta_g)=gW$.  The cone $\Xi$ is pointed: if both
$g$ and $-g$ lie in $\Xi$, the inequalities coming from
\[ e_4,e_5,e_6,e_7,e_1+e_4,e_2+e_5,e_3+e_7 \]
force all seven coordinates of $g$ to vanish.  Since
$\Mu\subseteq\Uneg\subseteq\UR$ by
\cref{prop:negative-fiber-theta-algebra,prop:upper-contained-in-invariants},
every triple-weight component of $\Mu$ is contained in the corresponding
finite-dimensional Kronecker weight space of $\UR$.  For the lowest-term
seed valuation fixed above, $\initer(\theta_g)=\mathbf x^g$, and its leaves
are one-dimensional.  Hence
\cref{lem:hilbert-theta-generation,lem:Xi-Hilbert-basis} shows that the
sixteen Hilbert-basis theta functions generate $\Mu$.

\begin{proposition}
\label{prop:twelve-generators-in-middle}
Write $f\doteq g$ when two nonzero homogeneous functions differ by a
scalar in $\kk^\times$.  For $g\in\mathcal H_1\cup\mathcal H_0$, the
table records the correspondence $\theta_g\doteq u_i$.
\begingroup
\setlength{\arraycolsep}{3pt}
\renewcommand{\arraystretch}{1.12}
\begin{equation}
\label{eq:theta-generator-table}
\begin{array}{cc|cc@{\qquad}cc}
\multicolumn{2}{c|}{\mathcal H_1}
 &\multicolumn{4}{c}{\mathcal H_0}\\
 g&\theta_g&g&\theta_g&g&\theta_g\\ \hline
 e_1&u_1&e_4&u_2&-e_1+e_2+e_4+e_6&u_{11}\\
 e_2&u_5&e_5&u_3&-e_2+e_3+e_5&u_6\\
 e_3&u_7&e_6&u_4&e_1-e_2+e_5+e_7&u_{12}\\
 2e_1-e_3+e_7&u_8&e_7&u_{13}&e_1-e_3+e_4+e_7&u_{10}
\end{array}
\end{equation}
\endgroup
In particular, every generator of $\UR$ except $u_9$ belongs to
$\Mu$.
\end{proposition}

\begin{proof}
For the degree-one labels and the four frozen basis vectors, the assertion
follows from their $g$-vectors.  The remaining four theta functions lie in
$\Mu\subseteq\Uneg\subseteq\UR$.  Their $W$-weights, in the order shown in
the last two columns of the table, are
\[ (321;33;42),\quad(210;21;21),\quad (321;42;33),\quad(221;32;32). \]
The corresponding invariant weight spaces are one-dimensional by
\cref{lem:relevant-one-dimensional-weights}; hence each nonzero theta function
is a scalar multiple of the indicated normalized invariant.  The scalars play
no role in algebra generation.
\end{proof}

Put
\[ \udisc:=u_9=u_{220,22}^{22}, \qquad \chi_\Delta=\wt(\udisc)=(2,2,0;2,2;2,2). \]

\begin{lemma}\label{lem:discriminant-chart-identities}
We have the following relations:
\begin{equation}
\label{eq:discriminant-chart-identities}
 \begin{aligned}
 u_1^2\udisc&=u_6^2+4u_2u_3u_4,\\
 u_5^2\udisc&=u_{10}^2-4u_2u_3u_{13}.
 \end{aligned}
\end{equation}
\end{lemma}

\begin{proof}
The first identity is the second relation in
\eqref{eq:basic-relations}.  The first and third relations there give
\[ u_1u_{10}=u_2u_8+u_3u_7, \qquad u_5u_6=u_3u_7-u_2u_8. \]
Hence, using also the last basic relation,
\[
 \begin{aligned}
 u_1^2\bigl(u_{10}^2-\udisc u_5^2\bigr)
 &=(u_2u_8+u_3u_7)^2-u_5^2(u_6^2+4u_2u_3u_4)\\
 &=4u_2u_3(u_7u_8-u_4u_5^2)
  =4u_1^2u_2u_3u_{13}.
 \end{aligned}
\]
Canceling $u_1^2$ in the domain $\UR$ proves the second identity.
\end{proof}

\begin{theorem}
\label{thm:middle-plus-discriminant}
The invariant ring, the four-chart intersection on the negative fiber, and
the one-generator extension of the middle algebra coincide:
\begin{equation}
\label{eq:upper-middle-extension}
 \UR=\Uneg=\Mu[\udisc].
\end{equation}
Moreover, $\udisc\notin\Mu$.
\end{theorem}

\begin{proof}
By \cref{prop:twelve-generators-in-middle}, the twelve generators
$u_i$ with $i\ne9$ belong to $\Mu\subseteq\Uneg$.  We claim that
$\udisc$ also belongs to all four Laurent charts.  In
$\cL_\circ^-$, $\cL_5^-$, and $\cL_7^-$, the element $u_1$ is
invertible, so the first identity in
\eqref{eq:discriminant-chart-identities} expresses $\udisc$ as a Laurent
function there.  In $\cL_1^-$, the element $u_5$ is invertible, and the
second identity does the same.  Thus $\udisc\in\Uneg$.

The thirteen-generator theorem and
\cref{prop:negative-fiber-theta-algebra,prop:upper-contained-in-invariants}
now give the closed chain
\[ \UR=\kk[u_1,\ldots,u_{13}] \subseteq\Mu[\udisc] \subseteq\Uneg \subseteq\UR, \]
which proves \eqref{eq:upper-middle-extension}.

For the final assertion, solve $gW=\chi_\Delta$.  Every solution has the form
\[ g=(2t-2,\,2s-2,\,2-2s-2t,\,s+t,\,2-s-t,\,s,\,t). \]
The inequalities $g_6,g_7\ge0$,
$g_2+g_3+g_7\ge0$,
$g_1+g_4\ge0$, and $g_3+g_7\ge0$, all belonging to
\eqref{eq:explicit-cover-cone}, force successively
$s,t\ge0$, $t=0$, $s\ge2$, and $s\le1$, a contradiction.
Hence there is no middle theta function of weight $\chi_\Delta$.  Since the
theta functions form a homogeneous basis, the corresponding middle weight
space is zero, whereas $\udisc\ne0$.
\end{proof}

\section{A polyhedral formula}
\label{sec:polyhedral}

\subsection{Paired theta generators}

Define the \emph{mutable degree} to be the $\ZZ$-grading determined by
\[ \deg u_1=\deg u_5=\deg u_7=1, \qquad \deg u_2=\deg u_3=\deg u_4=\deg u_{13}=0. \]
The three exchange relations are homogeneous for this assignment, so mutable
degree is a genuine $\ZZ$-grading of the cluster family, its negative
fiber, and $\Mu[\udisc]$.  In initial Laurent coordinates,
\[ \deg\theta_g=g_1+g_2+g_3, \qquad \deg\udisc=-2. \]
Order the positive- and negative-degree parts of the
Hilbert basis in \eqref{eq:Xi-Hilbert-basis} as
\[
\begin{aligned}
 h_1^+&=e_1,&
 h_1^-&=-e_1+e_4+e_5+e_6,&
 h_2^+&=e_2,&
 h_2^-&=-e_2+e_4+e_5+e_7,\\ 
 h_3^+&=e_3,&
 h_3^-&=-e_3+2e_4+e_6+e_7,&
 h_4^+&=2e_1-e_3+e_7,&
 h_4^-&=-2e_2+e_3+2e_5+e_7.
\end{aligned}
\]
By \eqref{eq:theta-generator-table},
\[ \theta_{h_1^+}\doteq u_1,\qquad \theta_{h_2^+}\doteq u_5,\qquad \theta_{h_3^+}\doteq u_7,\qquad \theta_{h_4^+}\doteq u_8. \]

\begin{proposition}\label{prop:paired-theta-generators}
For $1\le i\le4$, there is a scalar $c_i\in\kk^\times$ such that
\begin{equation}
\label{eq:paired-theta-generators}
 \theta_{h_i^-}=c_i\udisc\,\theta_{h_i^+}.
\end{equation}
Thus the four degree-minus-one generators are, up to normalization,
$\udisc u_1,\udisc u_5,\udisc u_7,\udisc u_8$.
\end{proposition}

\begin{proof}
Multiplication by the weight matrix gives
\[ h_i^-W=h_i^+W+\chi_\Delta. \]
More explicitly, the four common weights and the corresponding monomials are
\[
\begin{array}{c|c|c}
 i&h_i^-W&\text{monomial of this weight in }\kk[u_1,\ldots,u_{13}]\\ \hline
 1&(320;32;32)&\udisc u_1\\
 2&(331;43;43)&\udisc u_5\\
 3&(431;44;53)&\udisc u_7\\
 4&(431;53;44)&\udisc u_8.
\end{array}
\]
By \cref{lem:relevant-one-dimensional-weights}, each monomial
spans the corresponding one-dimensional weight space of $\UR$.  Both
$\theta_{h_i^-}$ and
$\udisc\theta_{h_i^+}$ are nonzero elements of that space, which proves
\eqref{eq:paired-theta-generators}.
\end{proof}

% \begin{corollary}
% \label{cor:negative-degree-factorization}
% If $f\in\Mu$ is homogeneous of mutable degree $d<0$, then there is a
% homogeneous $f^\vee\in\Mu$ of degree $-d$ such that
% \[ f=\udisc^{-d}f^\vee. \]
% \end{corollary}

% \begin{proof}
% By \cref{lem:hilbert-theta-generation,lem:Xi-Hilbert-basis}, the algebra
% $\Mu$ is generated by eight degree-zero elements, the four
% $\theta_{h_i^+}$, and the four $\theta_{h_i^-}$.  Express $f$ as a
% homogeneous polynomial in these generators.  In each monomial let $A$ and
% $B$ be the numbers of degree-one and degree-minus-one factors, counted with
% multiplicity.  Then $A-B=d<0$, so the monomial contains at least
% $B-A=-d$ negative-degree factors.  Rewrite any $-d$ of them by
% \eqref{eq:paired-theta-generators}, factor out $\udisc^{-d}$, and do this
% term by term.  The remaining polynomial is homogeneous of degree $-d$.
% \end{proof}

% \begin{lemma}
% \label{lem:multiplication-by-discriminant}
% Let $f\in\Mu$ be homogeneous of mutable degree $d$.  If
% $0\le r\le d$, then
% \[ \udisc^rf\in\Mu. \]
% In particular, if $g\in\Xi$ and $0\le r\le\deg g$, then
% $\udisc^r\theta_g\in\Mu$.
% \end{lemma}

% \begin{proof}
% Use the same homogeneous generating set.  In a monomial of degree $d$, let
% $A$ and $B$ count its degree-one and degree-minus-one factors.  Then
% $A-B=d$, so $A\ge d\ge r$.  Replace any $r$ degree-one factors by
% the corresponding degree-minus-one factors using
% \eqref{eq:paired-theta-generators}; apart from a nonzero scalar, this
% multiplies the monomial by $\udisc^r$.  Applying the replacement to every
% monomial in a homogeneous expression for $f$ proves the claim.
% \end{proof}

\begin{corollary}
\label{cor:multiplication-by-discriminant}
Let $f\in\Mu$ be homogeneous of mutable degree $d$.  If
$0\le r\le d$ or $d \le r \le 0$, then
\[ \udisc^rf\in\Mu. \]
\end{corollary}
\begin{proof}
By \cref{lem:hilbert-theta-generation,lem:Xi-Hilbert-basis}, the algebra
$\Mu$ is generated by eight degree-zero elements, the four
$\theta_{h_i^+}$, and the four $\theta_{h_i^-}$.  Express $f$ as a
homogeneous polynomial in these generators.  In each monomial let $N_+$ and
$N_-$ be the numbers of degree-one and degree-minus-one factors, counted with
multiplicity. Then $N_+ - N_-=d$. If $d \ge 0$, then $N_+ = N_- + d \ge d\ge r \ge 0$, so we may replace any $r$ degree-one factors by the corresponding degree-minus-one factors using
\eqref{eq:paired-theta-generators}.
If $d < 0$, then $N_- = N_+ - d \ge - d \ge -r \ge 0$, so we may replace any $-r$ degree-minus-one factors by
the corresponding degree-one factors using
\eqref{eq:paired-theta-generators}.
Apart from a nonzero scalar, these replacements
multiply the monomial by $\udisc^r$.  Applying the replacement to every
monomial in a homogeneous expression for $f$ proves the claim.
\end{proof}



\subsection{Reduction to one theta level}

Let
\[ \chi=(\lambda_1,\lambda_2,\lambda_3; \mu_1,\mu_2;\nu_1,\nu_2) \]
be an admissible triple weight.  Put
\begin{equation}
\label{eq:A-and-level}
 A=\lambda_1+2\lambda_2-\mu_1-\nu_1,
 \qquad
 \ell_0=\max\left\{0,\left\lfloor\frac{A+1}{2}\right\rfloor\right\}.
\end{equation}
A weight space with a nondominant triple weight is understood to be zero, and
the corresponding cone slice below is empty.  The proof will show, in
particular, that a nonzero coefficient forces
$\chi-\ell_0\chi_\Delta$ to be dominant.

\begin{theorem}
\label{thm:polyhedral-Kronecker}
The Kronecker coefficient of $\chi$ is the number of lattice points in a
single slice of the cone:
\begin{equation}
\label{eq:polyhedral-Kronecker-count}
 g_{\lambda,\mu,\nu}
 =\#\left\{g\in\Xi\cap\ZZ^7\ \middle|\
        gW=\chi-\ell_0\chi_\Delta\right\}.
\end{equation}
\end{theorem}

\begin{proof}
A direct calculation from \eqref{eq:weight-matrix-W} shows that every theta
function of weight $\chi-\ell\chi_\Delta$ has mutable degree
\begin{equation}
\label{eq:level-degree}
 d_\ell=2\ell-A.
\end{equation}
By \cref{lem:kronecker-weight-space,thm:middle-plus-discriminant}, the desired
Kronecker coefficient is the dimension of
\[ \Mu[\udisc]_\chi =\sum_{\ell\ge0}\udisc^\ell \Muwt{\chi-\ell\chi_\Delta}. \]
This sum is finite in a fixed triple weight, because
$\chi-\ell\chi_\Delta$ is nondominant for $\ell\gg0$.  Moreover, every
vector in $\Muwt{\chi-\ell\chi_\Delta}$ has the same mutable degree
$d_\ell$, since \eqref{eq:level-degree} expresses mutable degree as a
linear function of the triple weight on this fiber.  We show that every
summand is contained in the summand of level $\ell_0$.

% Suppose first that $\ell\ge\ell_0$, and put $r=\ell-\ell_0$.  Since
% $2\ell_0\ge A$,
% \[ d_\ell-r=\ell+\ell_0-A\ge0. \]
% Thus $0\le r\le d_\ell$, and
% \cref{lem:multiplication-by-discriminant} gives
% \[ \udisc^r\Muwt{\chi-\ell\chi_\Delta} \subseteq\Muwt{\chi-\ell_0\chi_\Delta}. \]
% Multiplication by $\udisc^{\ell_0}$ yields
% \begin{equation}
% \label{eq:high-level-reduction}
%  \udisc^\ell\Muwt{\chi-\ell\chi_\Delta}
%  \subseteq
%  \udisc^{\ell_0}\Muwt{\chi-\ell_0\chi_\Delta}.
% \end{equation}

% Now suppose $\ell<\ell_0$.  Then $d_\ell<0$.  For
% $f\in\Muwt{\chi-\ell\chi_\Delta}$,
% \cref{cor:negative-degree-factorization} gives
% \[ f=\udisc^{-d_\ell}f^\vee \qquad\text{for some}\qquad f^\vee\in\Muwt{\chi-(\ell-d_\ell)\chi_\Delta} \]
% of mutable degree $-d_\ell>0$.  Put $\ell^\vee=\ell-d_\ell=A-\ell$. Then
% $\udisc^\ell f=\udisc^{\ell^\vee}f^\vee$.
% The positive degree implies $\ell^\vee\ge\ell_0$, so
% \eqref{eq:high-level-reduction} at level $\ell^\vee$ places this element in
% the level-$\ell_0$ summand as well.

Suppose first that $\ell\ge\ell_0$. Put $r=\ell-\ell_0$. Since $2\ell_0\ge A$,
\[ d_\ell-r=\ell+\ell_0-A\ge0. \]
Thus $0\le r\le d_\ell$, and
\cref{cor:multiplication-by-discriminant} gives
\[ \udisc^r\Muwt{\chi-\ell\chi_\Delta} \subseteq\Muwt{\chi-\ell_0\chi_\Delta}. \]
Multiplication by $\udisc^{\ell_0}$ yields
\begin{equation}
\label{eq:high-level-reduction}
 \udisc^\ell\Muwt{\chi-\ell\chi_\Delta}
 \subseteq
 \udisc^{\ell_0}\Muwt{\chi-\ell_0\chi_\Delta}.
\end{equation}

Now suppose $\ell<\ell_0$. Then $\ell_0 \ge \ell + 1 \ge 1$. By \eqref{eq:A-and-level} we have $A \ge 2\ell_0 - 1$. Therefore, 
\[d_\ell = 2\ell - A \le 2(\ell_0 - 1) - A < 0.\]
Thus \cref{cor:multiplication-by-discriminant} gives
\[ \udisc^{d_\ell}\Muwt{\chi-\ell\chi_\Delta} \subseteq \Muwt{\chi-(\ell - d_\ell)\chi_\Delta}. \]
Multiplication by $\udisc^{\ell - d_\ell}$ yields
\[ \udisc^\ell\Muwt{\chi-\ell\chi_\Delta} \subseteq \udisc^{\ell - d_\ell} \Muwt{\chi-(\ell - d_\ell)\chi_\Delta}. \]
Put $\ell^\vee=\ell-d_\ell$. Then 
\[\ell^\vee = A-\ell \ge A - (\ell_0 - 1) \ge \ell_0,\] 
so \eqref{eq:high-level-reduction} at level $\ell^\vee$ proves that this summand is contained in
the level-$\ell_0$ summand as well.

It follows that
\[ \Mu[\udisc]_\chi =\udisc^{\ell_0}\Muwt{\chi-\ell_0\chi_\Delta}. \]
Multiplication by $\udisc^{\ell_0}$ is injective because the ambient
invariant ring is a domain.  The theta basis of $\Mu$ is homogeneous and
linearly independent, so the dimension of the last weight space is the
lattice-point count in \eqref{eq:polyhedral-Kronecker-count}.
\end{proof}

\subsection{A closed finite sum}

Assume from now on that the shifted triple is dominant; if it is not, the
coefficient is zero by \cref{thm:polyhedral-Kronecker}.  Set
\begin{equation}
\label{eq:shifted-weight}
 \bar\chi=\chi-\ell_0\chi_\Delta
 =(L_1,L_2,L_3;M_1,M_2;N_1,N_2).
\end{equation}
Thus
\[ L_1+L_2+L_3=M_1+M_2=N_1+N_2. \]
Interchanging the last two partitions if necessary, assume
\[ M_1-M_2\le N_1-N_2, \]
which is equivalent to $M_2\ge N_2$.  Put
\begin{equation}
\label{eq:closed-form-parameters}
 \begin{aligned}
 a&=L_1-L_2,& b&=L_2-L_3,& c&=L_3,\\
 \rho&=M_1-M_2,& h&=M_2-N_2,&d&=M_1+N_1-L_1-2L_2.
 \end{aligned}
\end{equation}
The integer $d$ is the mutable degree of the shifted fiber and is
nonnegative by the definition of $\ell_0$.

\begin{theorem}
\label{thm:closed-formula}
For $x\in\ZZ$, set $m_x=b-h-2x$, and
\begin{equation}
\label{eq:closed-sum-bounds}
 \begin{aligned}
 \mathsf L(x)&=\max\left\{
 0,\left\lceil\frac{a-d-x}{2}\right\rceil,N_2-2c-x
 \right\},\\
 \mathsf U(x)&=\min\{0,m_x\}
 +\min\left\{
 a,\left\lfloor\frac{a+h+x}{2}\right\rfloor,M_2-L_2+x
 \right\}.
 \end{aligned}
\end{equation}
Then
\begin{equation}
\label{eq:closed-Kronecker-sum}
 g_{\lambda,\mu,\nu}
 =\sum_{x=0}^{\min\{b,\rho\}}
   \pos{\mathsf U(x)-\mathsf L(x)+1}.
\end{equation}
Consequently, on the lattice of admissible triple weights, these Kronecker
coefficients are piecewise quasi-polynomial with quasiperiod dividing two.
\end{theorem}

\begin{proof}
We count the points in \eqref{eq:polyhedral-Kronecker-count}.  Put
$\delta=-g$, so the cone inequalities become
$\langle\alpha,\delta\rangle\le0$ for all
$\alpha\in A_4\cup A_5\cup A_6\cup A_7$, and solve
$-\delta W=\bar\chi$.  The integral solution set is an affine lattice of rank two.  Projection
to
\[ x=-\delta_5, \qquad y=-\delta_6 \]
is an affine lattice isomorphism onto $\ZZ^2$: solving the remaining five
equations over $\ZZ$ gives the inverse
\begin{equation*}
 \begin{aligned}
 \delta_1&=-a-b+h+2x+2y,&
 \delta_2&=a-d-2y,&
 \delta_3&=b-h-2x,\\
 \delta_4&=x-b,\qquad
 \delta_5=-x,&
 \delta_6&=-y,&\delta_7&=x+y+c-N_2.
 \end{aligned}
\end{equation*}
Substitution in the fourteen cone inequalities first gives
\[ 0\le x\le\min\{b,\rho\}. \]
For a fixed $x$, the nonredundant lower bounds on $y$ are
\[ y\ge0, \qquad y\ge\left\lceil\frac{a-d-x}{2}\right\rceil, \qquad y\ge N_2-2c-x, \]
whose maximum is $\mathsf L(x)$.

For the upper bounds, put
\[ C_x=\left\lfloor\frac{a+h+x}{2}\right\rfloor. \]
The six nonredundant inequalities are
\[ y\le a,\quad y\le C_x,\quad y\le M_2-L_2+x, \]
\[ y\le a+m_x,\quad y\le C_x+m_x, \quad y\le M_2-L_2+x+m_x. \]
Their minimum is $\mathsf U(x)$.  The remaining apparent bound
\[ y\le\left\lfloor\frac{a+b-x}{2}\right\rfloor \]
is redundant: before taking floors, its right-hand side is the arithmetic
mean of
\[ \frac{a+h+x}{2} \quad\text{and}\quad \frac{a+2b-h-3x}{2}. \]
Thus, for fixed $x$, the number of integral values of $y$ is
$\pos{\mathsf U(x)-\mathsf L(x)+1}$.  Summing over the allowed values of
$x$ proves \eqref{eq:closed-Kronecker-sum}.

Only division by two occurs in the endpoints.  Subdivide the admissible
parameter space so that the active affine branches in
\eqref{eq:closed-sum-bounds}, the order of their finitely many crossover
points of $x$, and the affine zeros of
$\mathsf U(x)-\mathsf L(x)+1$ are fixed.  On each resulting chamber both
the summation interval and the active positive-part branch are fixed.  The
sum is therefore a finite sum of affine functions over intervals with
parity-dependent endpoints, hence a quasi-polynomial with quasiperiod
dividing two.
\end{proof}

\begin{remark}
The quasiperiod bound is sharp, although it need not be minimal on every
chamber.  Indeed, for $n\ge0$ and $\lambda=\mu=\nu=(n,n)$, one has $A=n$ and
$\ell_0=\lceil n/2\rceil$.  If $n$ is even, the shifted weight is zero
and the slice consists of the origin; if $n$ is odd, the shifted weight is
nondominant.  Therefore
\[ g_{(n,n),(n,n),(n,n)}=\frac{1+(-1)^n}{2}. \]
Thus exact period two occurs on this ray, while on other chambers the parity
dependence may cancel and the quasi-polynomial may have period one.
\end{remark}

The \texttt{finite-sum} task in \path{Kron322.py} symbolically verifies the affine inverse and all fourteen substituted inequalities.
It also compares the formula with the symmetric-group character formula for
all relevant triples up to a user-specified size; the default bound is twelve.

\begin{corollary}[All $n\times2\times2$ formats]
\label{cor:closed-formula-22n}
Let $n\ge1$, and let $\lambda,\mu,\nu$ be partitions of a common integer
$N$ satisfying
$\ell(\lambda)\le n$ and $\ell(\mu),\ell(\nu)\le2$.
If $\ell(\lambda)>4$, then $g_{\lambda,\mu,\nu}=0$.  Otherwise pad
$\lambda$ to four parts and put $r=\lambda_4$.  If
$\mu_2<2r$ or $\nu_2<2r$, then again
$g_{\lambda,\mu,\nu}=0$.  In all remaining cases set
\[ \bar\lambda=(\lambda_1-r,\lambda_2-r,\lambda_3-r), \qquad \bar\mu=(\mu_1-2r,\mu_2-2r), \qquad \bar\nu=(\nu_1-2r,\nu_2-2r). \]
Apply \cref{eq:A-and-level,eq:shifted-weight,eq:closed-form-parameters,eq:closed-sum-bounds}
to $(\bar\lambda,\bar\mu,\bar\nu)$.  Then
\[ g_{\lambda,\mu,\nu} =\sum_{x=0}^{\min\{b,\rho\}} \pos{\mathsf U(x)-\mathsf L(x)+1}. \]
Thus the same explicit finite sum computes every Kronecker coefficient of
format $n\times2\times2$.
\end{corollary}

\begin{proof}
If $\ell(\lambda)>4$, then
$\Schur_\lambda(\kk^2\otimes\kk^2)=0$, so the coefficient vanishes.
Assume $\ell(\lambda)\le4$.  Proposition~\ref{prop:det-extension}, equivalently
\cref{cor:rectangular-translation}, with $b=c=2$ and
$r=\lambda_4$ gives zero when $\mu_2<2r$ or $\nu_2<2r$, and
otherwise gives
\[ g_{\lambda,\mu,\nu}=g_{\bar\lambda,\bar\mu,\bar\nu}. \]
The first partition on the right has length at most three, so
\cref{thm:closed-formula} applies.  For $n\ge4$,
\cref{cor:stable-first-factor} shows that no further dependence on the first
ambient dimension occurs.  For $n\le3$, enlarging the first vector space to
$\kk^3$ and padding by zeros leaves the same highest-weight multiplicity.
\end{proof}


\appendix
\raggedbottom

\section{The cluster algebra of the partial invariant ring}
\label{app:partial-ring}

We recall the part of the flagged-Kronecker construction needed in the main
text.  For positive integers $l_1,l_2,m$, let $K_{l_1,l_2}^m$ be the
flagged $m$-arrow Kronecker quiver
\[ -1\longrightarrow-2\longrightarrow\cdots\longrightarrow-l_1 \overset{m}{\Longrightarrow} l_2\longrightarrow l_2-1\longrightarrow\cdots\longrightarrow1 \]
with dimension vector $\beta(i)=|i|$.  Its semi-invariant ring is
\[ \SI_\beta(K_{l_1,l_2}^m) =\kk[\Rep_\beta(K_{l_1,l_2}^m)]^{\prod_i\SL_{\beta(i)}}. \]
Contracting the two equioriented flag arms gives a multigraded,
$\GL_m$-equivariant isomorphism
\begin{equation}
\label{eq:flagged-Kronecker-transfer}
 \SI_\beta(K_{l_1,l_2}^m)
 \cong
 \kk[\kk^{l_1}\otimes\kk^{l_2}\otimes\kk^m]^{U_{l_1}\times U_{l_2}}.
\end{equation}
One way to see \eqref{eq:flagged-Kronecker-transfer} is to tensor the
coordinate ring of the central representation space with the coordinate
rings of the two basic affine spaces $\GL_{l_i}/U_{l_i}$, and then take
$\GL_{l_1}\times\GL_{l_2}$-invariants.  Equivalently, successive Cauchy
decompositions along the arms retain exactly the highest-weight line at each
central vertex.

For $(l_1,l_2,m)=(3,2,2)$, the right-hand side is
$\SR$.  In the variable order
\eqref{eq:double-initial-cluster}, its extended exchange matrix is
\begin{equation*}
 B_s=
 \begin{bmatrix}
  0& 1&-1&-1& 1& 0& 1&-1\\
 -1& 0& 2& 1& 0&-1& 0& 0\\
  1&-2& 0& 0&-1& 1& 0& 0
 \end{bmatrix}.
\end{equation*}
The rows are indexed by $s_1,s_2,s_3$; the columns are ordered by
all eight variables in \eqref{eq:double-initial-cluster}, with the last five
columns frozen.  The results
of \cite{Fk1,Fk2} give
\[ \SR=\Asup. \]

Let $\epsilon_1,\ldots,\epsilon_8$ be the standard basis in the order
\eqref{eq:double-initial-cluster}.  The noninitial primitive rays of the
$g$-vector cone which are relevant here are
\begin{equation*}
 \begin{aligned}
 \gamma_1'&=\epsilon_1-\epsilon_3+\epsilon_6,&
 \gamma_2'&=-\epsilon_2+2\epsilon_3+\epsilon_4,\\
 \gamma_2^\circ&=-\epsilon_2+\epsilon_3+\epsilon_4,&
 \gamma_3'&=-\epsilon_1+\epsilon_2+\epsilon_5+\epsilon_7,\\
 \gamma_{31}&=-\epsilon_3+\epsilon_6+\epsilon_7,&
 \gamma_{32}&=-\epsilon_1+\epsilon_3+\epsilon_5+\epsilon_7,\\
 \gamma_{312}&=-\epsilon_1+\epsilon_5+2\epsilon_7.
 \end{aligned}
\end{equation*}
Together with the five frozen rays, these generate the integral
$g$-vector cone.  The labels used in the main text are fixed by the
following table.
\begin{table}[ht]
\centering
\small
\caption{The fourteen generators of $\SR$ used in the two-slice proof.}
\label{tab:partial-generators}
\begin{tabular}{ccll}
\toprule
$i$&$g(s_i)$&$\wt(s_i)$&$\partial(s_i)$\\ \midrule
1&$\epsilon_1$&$(110;20;11)$&0\\
2&$\epsilon_2$&$(100;10;01)$&$s_3$\\
3&$\epsilon_3$&$(100;10;10)$&0\\
4&$\epsilon_4$&$(110;11;02)$&$s_9$\\
5&$\epsilon_5$&$(110;11;20)$&0\\
6&$\epsilon_6$&$(200;11;11)$&0\\
7&$\epsilon_7$&$(111;21;12)$&$-s_8$\\
8&$\epsilon_8$&$(111;21;21)$&0\\ \midrule
9&$\gamma_2^\circ$&$(110;11;11)$&$2s_5$\\
10&$\gamma_3'$&$(211;22;22)$&$s_{11}$\\
11&$\gamma_{32}$&$(211;22;31)$&0\\
12&$\gamma_{31}$&$(211;22;13)$&$2s_{10}$\\
13&$\gamma_1'$&$(210;21;12)$&$2s_5s_2$\\
14&$\gamma_{312}$&$(222;33;33)$&0\\
\bottomrule
\end{tabular}
\end{table}

\begin{proposition}[{\cite[Proposition~9.5]{Fk1}}]
\label{prop:partial-minimal-generators}
The algebra $\SR$ is minimally generated by the eight initial
extended cluster variables and the five generic cluster characters with
$g$-vectors
\[ \gamma_2^\circ,\quad\gamma_3',\quad \gamma_{31},\quad\gamma_{32},\quad\gamma_{312}. \]
\end{proposition}

The fourteen generators used in \cref{subsec:invariant-upper} are obtained
from this minimal set by adjoining the one-step cluster variable with
$g$-vector $\gamma_1'$.  This redundant choice makes all root strings in
\eqref{eq:partial-on-s-generators} short and is the reason the two-slice proof has
the simple form there.

\section{A regular-sequence computation}
\label{app:regular-sequence-certificate}

\subsection{The single colon calculation}
\label{app:regular-sequence}

Let $S=\kk[U_1,\ldots,U_{13}]$, a Noetherian polynomial ring, and let
$J\subset S$ be generated by
\begin{equation}
\label{eq:J-six-relations}
 \begin{aligned}
 &U_3U_7-U_2U_8-U_5U_6,\\
 &U_1^2U_9-U_6^2-4U_2U_3U_4,\\
 &U_1U_{10}-2U_2U_8-U_5U_6,\\
 &U_1U_{11}+2U_2U_4U_5+U_6U_7,\\
 &U_1U_{12}-2U_3U_4U_5+U_6U_8,\\
 &U_1^2U_{13}-U_7U_8+U_4U_5^2.
 \end{aligned}
\end{equation}
Let $\varphi:S\to\Ugen$ send $U_i$ to $u_i$, and put
$I=\ker\varphi$.  The first slice in
\eqref{eq:invariantization-table} shows that, after inverting $U_1$, the
relations \eqref{eq:J-six-relations} give the complete presentation: the
last five relations eliminate $U_9,U_{10},U_{11},U_{12},U_{13}$, and the
first relation is the sole relation among the remaining eight generators.
Consequently, $I S_{U_1}=J S_{U_1}$.  Since
$S/I=\Ugen$ is a domain and $u_1\ne0$, the ideal $I$ is
$U_1$-saturated.  Contracting from $S_{U_1}$ therefore gives
\[ I=J S_{U_1}\cap S=J:U_1^\infty. \]
Because $S$ is Noetherian, the ascending union defining
$J:U_1^\infty$ stabilizes after finitely many colon operations.

For reproducibility, take graded reverse lexicographic order with the variable
order
\[ U_7>U_6>U_5>U_{13}>U_8>U_4>U_{10}>U_{11}> U_3>U_9>U_{12}>U_2>U_1. \]
The ancillary file
\texttt{Kron322\_regular\_sequence.m2} constructs
$I=J:U_1^\infty$, puts $K=I+(U_1)$, and verifies the exact assertion
\[ \operatorname{trim}(K:U_2)=K. \]
Equivalently,
\[ (I+(U_1)):U_2=I+(U_1). \]
The file prints the running Macaulay2 version and contains an assertion that causes the computation to terminate with an error unless this equality holds.  The independent exact calculation in the \texttt{regular-sequence} task of
\path{Kron322.py} performs both saturations by elimination over
$\mathbb Q$ and verifies the stronger equality
\[ K:U_2^\infty=K. \]
The generated reproducibility report records the two elimination steps and
the equality of the reduced Gr\"obner bases.  Since the ideals
and distinguished elements are defined over $\mathbb Q$, and every
characteristic-zero field extension of $\mathbb Q$ is flat, the colon
identity persists after scalar extension to $\kk$.

It follows that multiplication by $U_2$ is injective on
\[ S/(I,U_1)\cong\Ugen/u_1\Ugen. \]
Thus $u_1,u_2$ is a regular sequence in $\Ugen$, proving
\cref{lem:Ugen-regular-sequence}.

\section{The folded Donaldson--Thomas interpretation}
\label{app:folded-DT}

This appendix identifies the cluster transformation underlying the pairing in
\cref{prop:paired-theta-generators}.  It is not used in the proof of the
polyhedral or closed formulas.

Two consequences of \eqref{eq:basic-relations} will be used:
\begin{equation}
\label{eq:DT-two-identities}
 \begin{aligned}
 u_1^2u_2u_{13}+u_5^2u_2u_4+u_7^2u_3
   &=u_1u_7u_{10},\\
 u_{10}^2-\udisc u_5^2&=4u_2u_3u_{13}.
 \end{aligned}
\end{equation}
Indeed, the first and third relations in \eqref{eq:basic-relations} give
$u_1u_{10}=u_2u_8+u_3u_7$.  Multiplying by $u_7$ and using the last
relation proves the first identity.  For the second, use
$u_5u_6=u_3u_7-u_2u_8$, the second relation, and then the last relation:
\[
 \begin{aligned}
 u_1^2\bigl(u_{10}^2-\udisc u_5^2\bigr)
 &=(u_2u_8+u_3u_7)^2-u_5^2(u_6^2+4u_2u_3u_4)\\
 &=4u_2u_3(u_7u_8-u_4u_5^2)
  =4u_1^2u_2u_3u_{13}.
 \end{aligned}
\]
Cancelling $u_1^2$ gives the claim.

\begin{proposition}\label{prop:DT-coordinate-action}
There is a normalized folded cluster DT transformation $\Phi$ of the
negative fiber of \eqref{eq:zeta-family-exchanges} whose pullback satisfies
\begin{equation}
\label{eq:DT-coordinate-action}
 \Phi^*(u_i)=\udisc u_i\quad(i=1,5,7,8),
 \qquad
 \Phi^*(u_2,u_3,u_4,u_{13})=(u_2,u_3,u_4,u_{13}).
\end{equation}
\end{proposition}

\begin{proof}
Use the six-vertex unfolding with fold orbits
$\{1,4\},\{2,5\},\{3,6\}$.  In the mutable-row convention, all six rows
and the first six columns are mutable; the remaining columns correspond to
$u_2,u_3,u_4,u_{13},\zeta$:
\[
 \widetilde B_u=
 \begin{pmatrix}
 0&-1& 1& 0&-1& 1&-1& 1&-1& 0&0\\
 1& 0&-1& 1& 0&-1& 1&-1& 0& 1&1\\
 -1&1& 0&-1& 1& 0& 0& 0& 1&-1&0\\
 0&-1& 1& 0&-1& 1&-1& 1&-1& 0&0\\
 1& 0&-1& 1& 0&-1& 1&-1& 0& 1&1\\
 -1&1& 0&-1& 1& 0& 0& 0& 1&-1&0
 \end{pmatrix}.
\]
Zhou's DT word for this unfolding is the following (see \cite[Section~5.1]{Zhou}):
\[ \mathbf m=(1,3,2,4,6,5,1,6,4,3,2,5). \]
After mutation along $\mathbf m$ and terminal
relabeling in the order $(4,1,3,5,2,6)$, the first ten columns return to
their initial values and the $\zeta$-column changes sign.  Thus the terminal
coefficient is $\zeta^{-1}$; the ideal $(\zeta+1)$ is preserved, so the
transformation descends to the negative fiber.

Applying the same twelve exchange mutations to the variables, identifying
each fold orbit, and setting $\zeta=-1$, gives, before normalization,
\[ u_1\longmapsto u_1V, \qquad u_5\longmapsto-u_5V, \qquad u_7\longmapsto u_7V, \]
where
\begin{equation}
\label{eq:DT-factor-V}
 V=
 \frac{
 (u_1^2u_2u_{13}+u_5^2u_2u_4+u_7^2u_3)^2
 -4u_1^2u_7^2u_2u_3u_{13}}
 {u_1^2u_5^2u_7^2}.
\end{equation}
The sign in the second coordinate is removed by the coefficient-torus
normalization which multiplies the second and fifth unfolded mutable
variables by $\zeta$ and multiplies $u_4$ by $\zeta^{-2}$.  The
exponent vector of this rescaling lies in the right kernel of the first ten
columns of $\widetilde B_u$, so it is an automorphism of the seed pattern.
On $\zeta=-1$ it changes only the second folded sign.

Finally, \eqref{eq:DT-two-identities} transforms \eqref{eq:DT-factor-V} into
\[ V=\frac{u_{10}^2-4u_2u_3u_{13}}{u_5^2}=\udisc. \]
This proves the first three identities in \eqref{eq:DT-coordinate-action}.
Applying $\Phi^*$ to
$u_7u_8=u_4u_5^2+u_1^2u_{13}$ gives the identity for $u_8$.
\end{proof}

Every homogeneous Laurent rational function $f$ of mutable degree $d$
therefore satisfies
\[ \Phi^*(f)=\udisc^d f. \]
In particular, $\Phi^*(\udisc)=\udisc^{-1}$.  It fixes the eight
degree-zero generators of $\Mu$, and by
\cref{prop:paired-theta-generators} it exchanges, up to nonzero scalars, the
four paired generators of degrees $1$ and $-1$.  Hence it preserves
$\Mu$.  This is the cluster transformation associated with the weight-space pairing
used in the main proof.


\section{The \texorpdfstring{$D_5$}{D5}-parabolic origin of the four-frozen seed}
\label{app:D5-parabolic}

This appendix interprets the four-frozen seed
\[ \mathbf u^\circ=(u_1,u_5,u_7\,;u_2,u_3,u_4,u_{13}) \]
using the maximal parabolic at the trivalent node of type $D_5$.  The
appendix is not used elsewhere.

\subsection{The trivalent parabolic of
\texorpdfstring{$\operatorname{Spin}_{10}$}{Spin(10)}}

Let
\[ V=\kk^3,\qquad A=V^\vee,\qquad B=C=\kk^2, \qquad W=B\otimes C. \]
Alternating forms on $B$ and $C$ define a nondegenerate symmetric form
on $W$ by
\[ \langle b\otimes c,b'\otimes c'\rangle_W =\omega_B(b,b')\omega_C(c,c'). \]
Equip $E=A\oplus W\oplus A^\vee$ with the split quadratic form pairing
$A$ with $A^\vee$ and restricting to this form on $W$.  Then
$G=\operatorname{Spin}(E)$ is of type $D_5$.  Let $P_3\subset G$
stabilize the isotropic three-plane $A$.  Deleting the trivalent node
$\alpha_3$ leaves $A_2\sqcup A_1\sqcup A_1$, so, up to a finite
central quotient,
\[ (L_{P_3})_{\mathrm{der}} \simeq\SL(A)\times\SL(B)\times\SL(C). \]
The standard parabolic grading \cite[Chapter~3]{CapSlovak} has
\[ \mathfrak g_0\simeq\mathfrak{gl}(A)\oplus\mathfrak{so}(W),\qquad \mathfrak g_{-1}\simeq A^\vee\otimes W\simeq V\otimes B\otimes C,\qquad \mathfrak g_{-2}\simeq\bigwedge^2A^\vee. \]
It may also be viewed as a Vinberg $\theta$-representation
\cite{VinbergTheta}.  Indeed, if $\gamma:\mathbb G_m\to G$ is the
grading cocharacter and $\epsilon$ is a primitive fifth root of unity,
then $\theta=\operatorname{Ad}(\gamma(\epsilon))$ has fixed Lie algebra
$\mathfrak g_0$, while $\mathfrak g_{-1}$ is its
$\epsilon^{-1}$-eigenspace.  Thus, up to the identity component and a
finite central quotient, $(L_{P_3},\mathfrak g_{-1})$ is a
$\theta$-group representation in Vinberg's sense.
The block decomposition of $\mathfrak{so}(A\oplus W\oplus A^\vee)$ gives
these three summands directly; the dimension check is
\[ 3+12+(9+6)+12+3=45=\dim\mathfrak{so}_{10}. \]
After identifying $\SL(A)$ with $\SL(V)$ by duality and choosing
compatible Borels, the Levi maximal unipotent subgroup acts as
$U_3\times U_2\times U_2$.  Hence
\[ \UR=\kk[\mathfrak g_{-1}]^{U_L},\qquad U_L=U_3\times U_2\times U_2. \]
The Kronecker grading is therefore the Levi highest-weight grading of the
degree $-1$ component.

\subsection{Orthogonal Pl\"ucker and Gram covariants}

The homogeneous space $G/P_3$ is $\operatorname{OGr}(3,10)$.  If $J$
is a Gram matrix of the form on $W$, the big cell of $G/P_3$ is parametrized by
\[
 (X,S)\in\operatorname{Hom}(A,W)\oplus\bigwedge^2A^\vee,
 \qquad
 \operatorname{colspan}
 \begin{pmatrix}
 I_3\\ X\\ S-\frac12X^{\mathsf T}JX
 \end{pmatrix}.
\]
The tensor representation is the slice $S=0$.  Order the ten rows as the
three $A$-rows, four $W$-rows, and three $A^\vee$-rows, and write
$p_{ijk}$ for the corresponding Pl\"ucker minor.  With compatible bases
and Borels, direct substitution identifies the seven seed functions, up to
nonzero scalars, as
\[
\begin{aligned}
 u_1&\doteq p_{234},&u_5&\doteq p_{456},&u_7&\doteq p_{468},\\
 u_2&\doteq p_{346},&u_3&\doteq p_{345},&u_4&\doteq p_{238},&
 u_{13}&\doteq p_{8,9,10}.
\end{aligned}
\]
Here $\doteq$ denotes equality up to an element of $\kk^\times$.  Two of
these identifications follow immediately from the big-cell matrix:
\[ p_{238}=-\frac12\langle v_1,v_1\rangle_W,\qquad p_{8,9,10}=\left(-\frac12\right)^3\det(X^{\mathsf T}JX). \]
Thus $p_{238}\doteq D_1$ and $p_{8,9,10}\doteq D_3$.

Write $X=(v_1,v_2,v_3)$, with $v_i\in W$, and put
\[ D_r=\det\bigl(\langle v_i,v_j\rangle_W\bigr)_{1\le i,j\le r}. \]
Up to the same scalar convention,
\[ u_4=D_1,\qquad u_9=D_2,\qquad u_{13}=D_3. \]
The pair $u_2,u_3$ consists of the two highest components of
$v_1\wedge v_2$ under
\[ \bigwedge^2(B\otimes C) \simeq (\Sym^2B\otimes\bigwedge^2C) \oplus(\bigwedge^2B\otimes\Sym^2C). \]
If $N=*(v_1\wedge v_2\wedge v_3)\in W$, then $u_5$ is a
highest-weight coordinate of $N$, while $u_7$ is a highest-weight
component of $v_1\wedge N$.  The frozen variables $u_2,u_3$ are
highest-weight coordinates in the two irreducible
$\operatorname{Spin}_4$-summands above; $u_4=D_1$ and $u_{13}=D_3$
are the first and third principal Gram determinants.  The additional
generator in $\UR=\Mu[u_9]$ is the second principal Gram determinant
$D_2$.  Under these identifications, the relation
\[ u_7u_8=u_4u_5^2+u_1^2u_{13} \]
is the corresponding split-orthogonal Pl\"ucker--Gram identity.

\subsection{The \texorpdfstring{$D_5$}{D5}-weight lattice and Gale duality}

Let $\Lambda_{322}$ be the lattice of integer triples
$(\lambda;\mu;\nu)$ with
$\lvert\lambda\rvert=\lvert\mu\rvert=\lvert\nu\rvert=n$.  Realize
\[ Q(D_5)=\{z\in\ZZ^5\mid z_1+\cdots+z_5\equiv0\pmod2\} \]
with simple roots
$\alpha_1=e_1-e_2,\alpha_2=e_2-e_3,\alpha_3=e_3-e_4,
\alpha_4=e_4-e_5,\alpha_5=e_4+e_5$.  The map
\[ \Psi(\lambda;\mu;\nu) =(-\lambda_3,-\lambda_2,-\lambda_1, \mu_1+\nu_1-n,\nu_1-\mu_1) \]
is an isomorphism $\Lambda_{322}\simeq Q(D_5)$.  Indeed, for
$z\in Q(D_5)$, take
\[ n=-(z_1+z_2+z_3),\quad \lambda=(-z_3,-z_2,-z_1),\quad \mu_1=\frac{n+z_4-z_5}{2},\quad \nu_1=\frac{n+z_4+z_5}{2}, \]
and then $\mu_2=n-\mu_1$, $\nu_2=n-\nu_1$.

For
\[ (x_1,\ldots,x_7)=(u_1,u_5,u_7,u_2,u_3,u_4,u_{13}), \]
let $\chi_i=\Psi(\wt(x_i))$.  The rows of the following matrix are the
seven $D_5$-weights, while $B_u$ is obtained from $B_{u,\zeta}$ by
deleting the auxiliary $\zeta$-column:
\[
 \Omega_{D_5}=
 \begin{pmatrix}
 0&0&-1&1&0\\
 -1&-1&-1&1&0\\
 -1&-1&-2&1&1\\
 0&-1&-1&1&1\\
 0&-1&-1&1&-1\\
 0&0&-2&0&0\\
 -2&-2&-2&0&0
 \end{pmatrix},
 \qquad
 B_u=
 \begin{bmatrix}
 0&-2& 2&-1& 1&-1& 0\\
 2& 0&-2& 1&-1& 0& 1\\
 -2& 2& 0& 0& 0& 1&-1
 \end{bmatrix}.
\]
Let $\chi:\ZZ^7\to Q(D_5)$ send $e_i$ to $\chi_i$.

\begin{proposition}[Gale duality of the four-frozen seed]
\label{prop:D5-Gale-duality}
The exchange matrix is an integral Gale dual of the seven $D_5$-weights:
\[ B_u\Omega_{D_5}=0, \]
and there is an exact sequence
\[ 0\longrightarrow\ZZ \xrightarrow{\,1\mapsto(1,1,1)\,}\ZZ^3 \xrightarrow{\,B_u^{\mathsf T}\,}\ZZ^7 \xrightarrow{\,\chi\,}Q(D_5)\longrightarrow0. \]
\end{proposition}

\begin{proof}
The identities
\[
\begin{gathered}
 \alpha_1=-\chi_1-\chi_2+\chi_4+\chi_5,\quad
 \alpha_2=\chi_1-\chi_2+\chi_3-\chi_4,\quad
 \alpha_3=-\chi_1,\\
 \alpha_4=\chi_1+\chi_2-\chi_3,\quad
 \alpha_5=\chi_1+\chi_2-\chi_3+\chi_4-\chi_5
\end{gathered}
\]
show that $\chi$ is surjective.  Integral row reduction gives
\[ \ker\chi= \{(2s,2t,-2s-2t,s+t,-s-t,t,s)\mid s,t\in\ZZ\} =\operatorname{im}(B_u^{\mathsf T}). \]
Finally, $B_u^{\mathsf T}r=0$ forces $r_1=r_2=r_3$.
\end{proof}
The \texttt{d5} task in \path{Kron322.py} checks the
weight conversion, Smith normal form index, matrix product, ranks, integral kernel, and
Hermite normal form.

The three rows of $B_u$ give the homogeneity identities
\[ 2\chi_3+\chi_5=2\chi_2+\chi_4+\chi_6,\qquad 2\chi_3+\chi_5=2\chi_1+\chi_4+\chi_7, \qquad 2\chi_2+\chi_6=2\chi_1+\chi_7. \]

\begin{center}
\scriptsize
\renewcommand{\arraystretch}{1.08}
\begin{tabular}{@{}>{\raggedright\arraybackslash}p{0.28\textwidth}
                    >{\raggedright\arraybackslash}p{0.65\textwidth}@{}}
\toprule
\textbf{Cluster datum}&\textbf{Type-$D_5$ interpretation}\\
\midrule
$3\times2\times2$ tensor space
& The degree $-1$ component $\mathfrak g_{-1}$ of
  $P_3\subset\operatorname{Spin}_{10}$.\\
Triple $U$-invariant algebra $\UR$
& The Levi highest-weight algebra $\kk[\mathfrak g_{-1}]^{U_L}$.\\
Seven seed functions
& Restricted orthogonal Pl\"ucker covariants on
  $\operatorname{OGr}(3,10)$.\\
Frozen pair $u_2,u_3$
& Highest-weight coordinates in the two irreducible summands of
  $\bigwedge^2(B\otimes C)$.\\
Frozen variables $u_4,u_{13}$
& The first and third principal Gram determinants $D_1,D_3$;
  $u_9=D_2$ is the second one.\\
Exchange data
& The rows of $B_u$ are the full lattice of $D_5$-weight relations;
  the third exchange is the Pl\"ucker--Gram identity.\\
Mutable Markov quiver
& The mutation dynamics of the chart; it is of infinite cluster type, not
  finite Dynkin type $D_5$.\\
\bottomrule
\end{tabular}
\end{center}

In this interpretation, type $D_5$ supplies the ambient representation, the
rank-five weight lattice, and the homogeneity relations; the doubled Markov
triangle supplies the mutation dynamics.

\section{Two further models}
\label{app:other-models}

For comparison, we record two additional coefficient-fiber models; neither is
used in the main text.  Here the invariants $u_1,\ldots,u_{13}$ retain the
normalization of \eqref{eq:basic-relations}, and the primed variables are
normalized by the formulas below.  Thus the scalar coefficients in the
exchange relations are fixed.  As in \cref{sec:slice-volume}, a negative
exchange sign denotes the fiber $\zeta=-1$ of an ordinary cluster family.
The matrices suppress the auxiliary $\zeta$-column and retain only the
mutable and invariant frozen columns.

\subsection{A four-mutable model}

Take mutable variables $(u_1,u_2,u_3,u_6)$ and frozen variables
$(u_7,u_8,u_9)$.  The exchange relations are
\begin{align*}
 u_1u_1'&=u_2u_8+u_3u_7,
 &u_2u_2'&=u_1^2u_9-u_6^2,\\
 u_3u_3'&=u_1^2u_9-u_6^2,
 &u_6u_6'&=u_3u_7-u_2u_8,
\end{align*}
where
\[ u_1'=u_{10},\qquad u_2'=4u_3u_4,\qquad u_3'=4u_2u_4,\qquad u_6'=u_5. \]
These formulas follow from the first three identities in
\eqref{eq:basic-relations}.  In the mutable-row convention, with columns ordered
by $(u_1,u_2,u_3,u_6;u_7,u_8,u_9)$, the exchange matrix is
\[
 B_A=
 \begin{bmatrix}
 0&-1& 1& 0& 1&-1& 0\\
 2& 0& 0&-2& 0& 0& 1\\
-2& 0& 0& 2& 0& 0&-1\\
 0& 1&-1& 0&-1& 1& 0
 \end{bmatrix}.
\]
In this case, 
\[ \UR=\mathcal M_A[u_4,u_{13}]. \]

\subsection{A five-mutable model}

Take mutable variables $(u_1,u_2,u_3,u_5,u_6)$ and frozen variables
$(u_{11},u_{12})$.  With the normalizations
\begin{equation*}
 \begin{aligned}
 u_1'&=u_5u_9,&
 u_2'&=-2u_3u_4u_5-u_6u_8,&
 u_3'&=2u_2u_4u_5-u_6u_7,\\
 u_5'&=u_1u_9,&
 u_6'&=-u_{10},
 \end{aligned}
\end{equation*}
\eqref{eq:basic-relations} gives the exchange relations:
\begin{align*}
 u_1u_1'&=u_2u_{12}-u_3u_{11},
 &u_2u_2'&=u_1u_3u_{11}+u_5u_6^2,\\
 u_3u_3'&=u_1u_2u_{12}-u_5u_6^2,
 &u_5u_5'&=u_2u_{12}-u_3u_{11},\\
 u_6u_6'&=u_2u_{12}+u_3u_{11}.
\end{align*}
With columns ordered by
$(u_1,u_2,u_3,u_5,u_6;u_{11},u_{12})$, the exchange matrix is
\[
 B_B=
 \begin{bmatrix}
 0& 1&-1& 0& 0&-1& 1\\
-1& 0&-1& 1& 2&-1& 0\\
 1& 1& 0&-1&-2& 0& 1\\
 0&-1& 1& 0& 0& 1&-1\\
 0&-1& 1& 0& 0& 1&-1
 \end{bmatrix}.
\]
In this case, 
\[ \UR=\mathcal M_B[u_4,u_9,u_{13}]. \]

For each model, the corresponding cone and paired theta generators give
analogues of \cref{thm:polyhedral-Kronecker,thm:closed-formula}: every
Kronecker coefficient is reduced to one theta level and then expressed as a
nonnegative finite lattice-point sum.  We omit the parallel formulas because
the four-frozen model gives the shortest presentation.

\bibliographystyle{amsalpha}
\bibliography{Kron322}

\end{document}
